\documentclass[a4paper,12pt]{article}

% ------------------------------------------------------------------------------
% 1. KONFIGURATION & SPRACH-LOGIK (Muss zuerst geladen werden)
% ------------------------------------------------------------------------------
\def\myType{2}
%% 0=Projektarbeit   %%
%% 1=Bachelorarbeit  %%
%% 2=sonstige Arbeiten   %%

% Titel laut Aufgabenstellung
\def\myTopic{Sudoku, SAT und Lösbarkeit}
\def\mySubTopic{Entwicklung eines Problemlösers}

\def\myStudyProgram{ Data Science und Künstliche Intelligenz} % Studiengang
\def\myNeedForFieldOfStudy{0} % Gibt an ob eine Studienrichtung benötigt wird
%% 0=keine Studienrichtung %%
%% 1=Studienrichtung %%
\def\myFieldOfStudy{}         % Studiengangsrichtung

% Auf dem Titelblatt untereinander, in den PDF-Metadaten durch Komma getrennt
\def\myAutor{\texorpdfstring{\begin{tabular}[t]{@{}l@{}}Raphael Henninger\\Nils Mauritz\\Devin Rohrer\\Diana Gast\end{tabular}}{Raphael Henninger, Nils Mauritz, Devin Rohrer, Diana Gast}}
\def\myCourse{WDS125}               
\def\myProf{Prof. Dr. Aykan Inan}
\def\myCompany{}
\def\myCompanion{}
\def\myEndDate{22.11.2026}              


%% Projektarbeit %%
\def\myProjNumber{1}          % [1|2]

%% Bachelorarbeit %%
\def\myTypeOfBachelor{}       % [Arts|Science]

%% sonstige Arbeiten %%
\def\myTypeOfWork{Projektdokumentation} % TODO: Bezeichnung laut Aufgabenstellung prüfen
\def\myModule{Theoretische Informatik}  % TODO: genauen Modulnamen prüfen
\def\myNames{1}               % Gibt an ob mehrere Personen an der Arbeit mitwirken
%% 0=Einzahl %%
%% 1=Mehrzahl %%

%% Konfiguration der Verzeichnisse %%
% für alle folgenden Konfigurationen gilt: 0=kein Verzeichnis | 1=Verzeichnis %

\def\myAbrevationDirectory{1}   % Abkürzungsverzeichnis
\def\myIllustrationDirectory{1} % Abbildungsverzeichnis
\def\myTableDirectory{1}        % Tabellenverzeichnis
\def\myListingDirectory{1}      % Listingverzeichnis
\def\myFormularDirectory{0}     % Formelverzeichnis
\def\mySymbolDirectory{0}       % Symbolverzeichnis
\def\myKiDirectory{1}           % Verzeichnis der verwendeten KI-Hilfsmittel

%% Konfiguration Selbständigkeitserklärung und Sperrvermerk %%
\def\myDecOfIndependence{1}
%% 0=keine Selbständigkeitserklärung   %%
%% 1=Selbständigkeitserklärung    %%

\def\myBlockingNote{0}
%% 0=kein Sperrvermerk   %%
%% 1=Sperrvermerk    %%

% Ort und Datum der Unterschrift auf Selbständigkeitserklärung und Sperrvermerk
\def\myDate{22.11.2026}                 % TODO
\def\myPlace{Ravensburg}

% Sprachauswahl: 'ngerman' oder 'english'
\newcommand{\myLanguage}{ngerman}
           % Lädt \myLanguage und andere Einstellungen
\usepackage{ifthen}             % Ermöglicht If-Abfragen für die Sprache

% Sprache laden (basierend auf \myLanguage aus config.tex)
% WICHTIG: inputenc sollte davor stehen
\usepackage[utf8]{inputenc}
\usepackage[\myLanguage]{babel} 

% Sprach-Definitionen laden (Deutsch oder Englisch)
\ifthenelse{\equal{\myLanguage}{ngerman}}
    {% Datei: main.tex
\newcommand{\langAnhang}{Anhang}       
\newcommand{\langFigureName}{Abb.}     
\newcommand{\langTableName}{Tab.}      

% Datei: 10_titel_bachelor.tex
\newcommand{\langBachelorThesis}{Bachelorarbeit}
\newcommand{\langForDegree}{für die\\Prüfung zum Bachelor of}
\newcommand{\langFaculty}{an der Fakultät Wirtschaft}
\newcommand{\langStudyProgram}{im Studiengang}
\newcommand{\langFieldOfStudy}{in der Fachrichtung}
\newcommand{\langAt}{an der}
\newcommand{\langAutor}{Autor/-in}
\newcommand{\langKurs}{Kurs}
\newcommand{\langBetreuer}{DHBW-Betreuer/-in}
\newcommand{\langFirma}{Dualer Partner}
\newcommand{\langBegleiter}{Begleiter/-in Dualer Partner}
\newcommand{\langAbgabefrist}{Abgabefrist}

% Datei: 10_titel_projekt_doku.tex
\newcommand{\langAutoren}{Autoren/Autorinnen}
\newcommand{\langModul}{Modul}
\newcommand{\langDozent}{Dozent/-in}

% Datei: 10_titel_projekt.tex
\newcommand{\langProjektarbeit}{Projektarbeit}

% Datei: 13_abkuerzungsverzeichnis.tex
\newcommand{\langAbkuerzungsverzeichnis}{Abkürzungsverzeichnis}

% Datei: 14_Formelverzeichnis.tex
\newcommand{\langFormelverzeichnis}{Formelverzeichnis}
\newcommand{\langFormelPrefix}{FV.}

% Datei: 17_listingverzeichnis.tex
\newcommand{\langListingsverzeichnis}{Listingverzeichnis}
\newcommand{\langListingName}{Listing}

% Datei: 18_literaturverzeichnis.tex
\newcommand{\langLiteraturverzeichnis}{Quellen- und Literaturverzeichnis}

% Datei: 19_symbolverzeichnis.tex
\newcommand{\langSymbolverzeichnis}{Symbolverzeichnis}
} % Lädt deutsche Texte & Anpassungen
    {% Datei: main.tex
\newcommand{\langAnhang}{Appendix}     
\newcommand{\langFigureName}{Fig.}     
\newcommand{\langTableName}{Tab.}     

% Datei: 10_titel_bachelor.tex
\newcommand{\langBachelorThesis}{Bachelor Thesis}
\newcommand{\langForDegree}{for the\\examination for the degree of Bachelor of}
\newcommand{\langFaculty}{at the School of Business}
\newcommand{\langStudyProgram}{in the study program}
\newcommand{\langFieldOfStudy}{in the field of study}
\newcommand{\langAt}{at the}
\newcommand{\langAutor}{Author}
\newcommand{\langKurs}{Course}
\newcommand{\langBetreuer}{DHBW Supervisor}
\newcommand{\langFirma}{Dual Partner}
\newcommand{\langBegleiter}{Company Supervisor}
\newcommand{\langAbgabefrist}{Submission Date}

% Datei: 10_titel_projekt_doku.tex
\newcommand{\langAutoren}{Authors}
\newcommand{\langModul}{Module}
\newcommand{\langDozent}{Lecturer}

% Datei: 10_titel_projekt.tex
\newcommand{\langProjektarbeit}{Project Work}

% Datei: 13_abkuerzungsverzeichnis.tex
\newcommand{\langAbkuerzungsverzeichnis}{List of Abbreviations}

% Datei: 14_Formelverzeichnis.tex
\newcommand{\langFormelverzeichnis}{List of Formulas}
\newcommand{\langFormelPrefix}{LF.}

% Datei: 17_listingverzeichnis.tex
\newcommand{\langListingsverzeichnis}{List of Listings}
\newcommand{\langListingName}{Listing}

% Datei: 18_literaturverzeichnis.tex
\newcommand{\langLiteraturverzeichnis}{Bibliography}

% Datei: 19_symbolverzeichnis.tex
\newcommand{\langSymbolverzeichnis}{List of Symbols}} % Lädt englische Texte & Anpassungen

% ------------------------------------------------------------------------------
% 2. LAYOUT & STANDARD-PAKETE
% ------------------------------------------------------------------------------
\usepackage[left=2.5cm,right=2.5cm,top=2.5cm,bottom=2cm,includehead]{geometry}
\usepackage[hyperfootnotes=false, pdfborder={0 0 1}]{hyperref} 
\usepackage[nottoc]{tocbibind}
\usepackage{makeidx}
\usepackage[intoc]{nomencl}
\usepackage{fancyhdr}
\usepackage{pdfpages}
\usepackage{amsmath}
\usepackage[labelfont=bf, aboveskip=1mm, font={small, singlespacing}]{caption}
\usepackage{setspace}
\usepackage[bottom,multiple,hang,marginal]{footmisc}
\usepackage{wrapfig}
\usepackage{nomencl} 
\usepackage{chngcntr}
\usepackage{import}
\usepackage{graphicx}
\usepackage{tabularx}
\usepackage{longtable}
\usepackage{color}
\usepackage{enumitem}
\usepackage{listings}
\usepackage[noend]{algpseudocode}
\algrenewcommand\algorithmicrequire{\textbf{Eingabe:}}
\algrenewcommand\algorithmicensure{\textbf{Ausgabe:}}
\algrenewcommand\algorithmicfunction{\textbf{Funktion}}
\algrenewcommand\algorithmicif{\textbf{falls}}
\algrenewcommand\algorithmicthen{\textbf{dann}}
\algrenewcommand\algorithmicelse{\textbf{sonst}}
\algrenewcommand\algorithmicwhile{\textbf{solange}}
\algrenewcommand\algorithmicforall{\textbf{für alle}}
\algrenewcommand\algorithmicdo{\textbf{tue}}
\algrenewcommand\algorithmicrepeat{\textbf{wiederhole}}
\algrenewcommand\algorithmicuntil{\textbf{bis}}
\algrenewcommand\algorithmicreturn{\textbf{Rückgabe}}
\usepackage{zref}
\usepackage{wasysym}
\usepackage{amssymb}
\usepackage{xspace}
\usepackage{multicol}
\usepackage{booktabs}
\usepackage{siunitx}
\RequirePackage{multirow}
\RequirePackage{mwe}
\RequirePackage{etoolbox}
\usepackage{float}
\usepackage{adjustbox}
\usepackage{comment}
\usepackage{subcaption}
\usepackage{acronym}
\usepackage{csquotes}
\usepackage{forest}                % Baumdiagramme (lädt TikZ), z. B. Suchbaum im Kapitel Backtracking
\usepackage{algorithm}             % Gleitumgebung für Pseudocode
\usepackage[noend]{algpseudocode}  % Pseudocode-Befehle (\Function, \If, \For, ...)
\floatname{algorithm}{Algorithmus}
\numberwithin{algorithm}{section}
\AtBeginDocument{\counterwithin{lstlisting}{section}} % Zähler wird von listings erst bei \begin{document} angelegt
\makeatletter
\def\ext@algorithm{lol}                            % Algorithmen ins gemeinsame Listingverzeichnis schreiben (Dokumentreihenfolge)
\def\l@algorithm{\@dottedtocline{1}{1.5em}{2.3em}} % gleiche Formatierung wie Listing-Einträge
\makeatother

\pdfminorversion=6

%%%%%%%%%%%%%%%%%%%%%%%% Eigene Farbwerte definieren %%%%%%%%%%%%%%%%%%%%%%%%
\definecolor{boxgray}{gray}{0.9}     
\definecolor{commentgray}{gray}{0.5} 
\definecolor{darkgreen}{rgb}{0,.5,0} 

%%%%%%%%%%%%%%%%%%%%%%%% Eigene Kommandos definieren %%%%%%%%%%%%%%%%%%%%%%%%
\newcommand{\gqq}[1]{\glqq #1\grqq}

\providecommand*{\footref}[1]{
	\begingroup
		\unrestored@protected@xdef\@thefnmark{\ref{#1}}
	\endgroup
\@footnotemark}

\newcommand{\mypageref}[1]{\ref{#1} \nameref{#1} auf Seite \pageref{#1}}

\newcommand{\myboxquote}[1]{
	\begin{quotation}
		\colorbox{boxgray}{\parbox{0.78\textwidth}{#1}}
	\end{quotation}
	\vspace*{1mm}
}

% Definition von \vgl{#1}{#2} (APA-konform)
\newcommand{\vgl}[2][]{\parencite[vgl.][#1]{#2}}


\makeatletter
\zref@newprop*{appsec}{}
\zref@addprop{main}{appsec}

\def\applabel#1#2{%
	\zref@setcurrent{appsec}{#2}%
	\zref@wrapper@immediate{\zref@label{#1}}%
}

\def\appref#1{%
	\hyperref[#1]{\zref@extract{#1}{appsec}}%
}
\makeatother

% Definition von \appsection: Nutzt jetzt \langAnhang Variable
\newcommand{\theappsection}[1]{\langAnhang \Alph{section}:~\protect #1}
\newcommand{\appsection}[2]{
	\addtocounter{section}{1}
	\phantomsection
	\addcontentsline{toc}{section}{\theappsection{#1}}
	\markboth{\theappsection{#1}}{}

	\section*{\theappsection{#1}}
	\applabel{#2}{\langAnhang \Alph{section}}
	\label{#2}
}

% Nummerierung der Formeln
\renewcommand{\theequation}{\arabic{equation}}

%%%%%%%%%%%%% Index, Abkürzungsverzeichnis %%%%%%%%%%%%
\makeindex
\makenomenclature

% Zitation (APA 7 mit DHBW Sonderregel maxcitenames=3)
\usepackage[style=apa, backend=biber, language=ngerman, maxcitenames=3, mincitenames=1]{biblatex}
\addbibresource{literatur.bib}

%%%%%%%%%%%%%%%%%%%%%%%%%%%%%%% PDF-Optionen %%%%%%%%%%%%%%%%%%%%%%%%%%%%%%%%
\hypersetup{
	bookmarksopen=false,
	bookmarksnumbered=true,
	bookmarksopenlevel=0,
	pdftitle=\myTopic,
	pdfsubject=\myTopic,
	pdfauthor=\myAutor,
	pdfborder={0 0 1}, 
	pdfstartview=Fit,
	pdfpagelayout=SinglePage
}

%%%%%%%%%%%%%%%%%%%%%%%%%%%% Kopf- und Fußzeile %%%%%%%%%%%%%%%%%%%%%%%%%%%%%
\pagestyle{fancy}
\fancyhf{} 
\fancyfoot[R]{\thepage}
\renewcommand{\headrulewidth}{0.5pt} 
\renewcommand{\footrulewidth}{0pt} 

%%%%%%%%%%%%%%%%%%%%%%%%% Layout und Beschriftungen %%%%%%%%%%%%%%%%%%%%%%%%%
\onehalfspacing
\setlist{noitemsep}

\setlength{\parskip}{6pt}
\setlength{\parindent}{0pt}

% Dynamische Anpassung der Bild/Tabellen-Unterschriften basierend auf der Sprache
\addto\captionsngerman{
  \renewcommand{\figurename}{\langFigureName}
  \renewcommand{\tablename}{\langTableName}
}
\addto\captionsenglish{
  \renewcommand{\figurename}{\langFigureName}
  \renewcommand{\tablename}{\langTableName}
}

\numberwithin{table}{section}
\numberwithin{figure}{section}
\renewcommand{\thetable}{\arabic{section}.\arabic{table}}
\renewcommand{\thefigure}{\arabic{section}.\arabic{figure}}
\renewcommand{\thefootnote}{\arabic{footnote}}
\renewcommand{\multfootsep}{, }

%%%%%%%%%%%%%%%%%%%%%%%%%%%%%%% Listingstyle %%%%%%%%%%%%%%%%%%%%%%%%%%%%%%%%
\lstset{
	basicstyle=\ttfamily\scriptsize,
	commentstyle=\color{commentgray}\textit,
	showstringspaces=false,
	stringstyle=\color{darkgreen},
	keywordstyle=\color{blue},
	numbers=left,
	numberstyle=\tiny,
	stepnumber=1,
	numbersep=15pt,
	tabsize=2,
	fontadjust=true,
	frame=single,
	backgroundcolor=\color{boxgray},
	captionpos=b,
	linewidth=0.94\linewidth,
	xleftmargin=0.1\linewidth,
	breaklines=true,
	aboveskip=16pt
}

%%%%%%%%%%%%%%%%%%%%%%%%%%%%%%%%%%%%%%%%%%%%%%%%%%%%%%%%%%%%%%%%%%%%%%%%%%%%%
%%                         DOKUMENT BEGINN                                 %%
%%%%%%%%%%%%%%%%%%%%%%%%%%%%%%%%%%%%%%%%%%%%%%%%%%%%%%%%%%%%%%%%%%%%%%%%%%%%%

\begin{document}
	\pagenumbering{Roman}
	
	% Titelblatt (Dateien bleiben gleich, Variablen werden innen genutzt)
	\ifcase\myType
            \begin{titlepage}
	\begin{center}
		\vspace*{1cm}
		\LARGE\bf\myTopic\\
		\Large\rm\mySubTopic\\
		\vspace*{2cm}
		\bf \myProjNumber.~\langProjektarbeit\\
		\vspace*{0.5cm}\singlespacing
		\normalsize\rm
		\ifcase\myNeedForFieldOfStudy
            \langFaculty\\
            \langStudyProgram \myStudyProgram \\
        \or 
            \langFaculty\\
            \langStudyProgram \myStudyProgram \\
            \langFieldOfStudy \myFieldOfStudy\\
        \else
        \fi
		\vspace*{0.5cm}\singlespacing
		\langAt\\
		DHBW Ravensburg
		\vfill
	\end{center}
	\begin{tabular}{lll}
		\langAutor: &\myAutor\\
            \langKurs: &\myCourse\\
            \langBetreuer: &\myProf\\
		\langFirma: &\myCompany\\
		\langBegleiter: &\myCompanion \\
		\langAbgabefrist: &\myEndDate
	\end{tabular}
	\newline
	\vspace*{1cm}
	\newline
	\begin{tabularx}{\textwidth}{l@{\extracolsep\fill}r}
	  % Unterschrift des verantwortlichen Ausbilders&\\
	 % (oder des Personalverantwortlichen)&\rule{6cm}{0.3mm}\\
	\end{tabularx}
\end{titlepage}
\newpage
        \or
            \begin{titlepage}
	\begin{center}
		\vspace*{2cm}
		\LARGE\bf\myTopic\\
		\Large\rm\mySubTopic\\
		\vspace*{3cm}
		\bf \langBachelorThesis\\
		\normalsize\rm
		\vspace*{0.5cm}\singlespacing
		\langForDegree \myTypeOfBachelor\\
		\vspace*{0.5cm}\singlespacing
        \ifcase\myNeedForFieldOfStudy
        \langFaculty\\
        \langStudyProgram \myStudyProgram \\
        \or
            \langFaculty\\
        \langStudyProgram \myStudyProgram \\
            \langFieldOfStudy \myFieldOfStudy\\
        \else
        \fi
		\vspace*{0.5cm}\singlespacing
		\langAt\\
		DHBW Ravensburg
		\vfill
	\end{center}
	\begin{tabular}{ll}
		\langAutor: &\myAutor\\
        \langKurs: &\myCourse\\
        \langBetreuer: &\myProf\\
		\langFirma: &\myCompany\\
		\langBegleiter: &\myCompanion \\
		\langAbgabefrist: &\myEndDate
	\end{tabular}
\end{titlepage}
\newpage
        \or
            \begin{titlepage}
	\begin{center}
		\vspace*{1cm}
		\LARGE\bf\myTopic\\
		\Large\rm\mySubTopic\\
		\vspace*{2cm}
		\bf \myTypeOfWork \\
		\vspace*{0.5cm}\singlespacing
		\normalsize\rm
		\ifcase\myNeedForFieldOfStudy
            \langFaculty\\
            \langStudyProgram \myStudyProgram \\
        \or 
            \langFaculty\\
            \langStudyProgram \myStudyProgram \\
            \langFieldOfStudy \myFieldOfStudy\\
        \else
        \fi
		\vspace*{0.5cm}\singlespacing
		\langAt\\
		DHBW Ravensburg
		\vfill
	\end{center}
	\begin{tabular}{ll}
		\ifcase\myNames
            \langAutor:
        \or 
            \langAutoren:
        \else
        \fi
        &\myAutor\\
        \langModul: &\myModule\\
		\langDozent:&\myProf\\
		\langAbgabefrist:&\myEndDate
	\end{tabular}
	
\end{titlepage}
\newpage
        \else
        \fi

    \ifcase\myBlockingNote
        \or
            \thispagestyle{empty} %change

\begin{center}
	\vspace*{1cm}
	\Huge\bf Sperrvermerk\\
	\vspace*{2cm}
	\large\rm
	
	\begin{quotation}
		\hspace*{-1.75em}
		\parbox{0.85\textwidth}{\singlespacing Der Inhalt dieser Arbeit darf weder als Ganzes noch in Auszügen Personen außerhalb des Prüfungsprozesses und des Evaluationsverfahrens zugänglich gemacht werden, sofern keine anderslautende Genehmigung des Dualen Partners vorliegt.}
	\end{quotation}
	\vspace*{0.5cm}
	\begin{quotation}
		\hspace*{-1.5em}
		\parbox{\textwidth}{
			\singlespacing
			\begin{tabularx}{0.83\textwidth}{l@{\extracolsep\fill}l}
				\rule{4cm}{0.3mm}&\rule{4cm}{0.3mm}\\
				\large Ort, Datum&\large Unterschrift
			\end{tabularx}}
	\end{quotation}
	\end{center}
\newpage

        \else
    \fi

	\pagestyle{fancy}
    \setcounter{page}{1} % Reset für Inhaltsverzeichnis (Beginn bei I)
    
    % Header-Definition VOR die Metaseiten schieben
	\fancyhead[L]{\nouppercase{\leftmark}} 
	\tableofcontents
\newpage


\ifcase\myAbrevationDirectory

\or
    \newpage
\section*{\langAbkuerzungsverzeichnis}

\markboth{\langAbkuerzungsverzeichnis}{\langAbkuerzungsverzeichnis} 

\addcontentsline{toc}{section}{\langAbkuerzungsverzeichnis}

\begin{acronym}[DHBW]
    \acro{ALO}{at least once}
    \acro{AMO}{at most once}
    \acro{CDCL}{Conflict-Driven Clause Learning}
    \acro{DPLL}{Davis-Putnam-Logemann-Loveland-Algorithmus}
    \acro{FC}{Forward Checking}
    \acro{KNF}{Konjunktive Normalform}
    \acro{LBD}{literal block distance}
    \acro{MRV}{Minimum Remaining Values}
    \acro{SAT}{Erfüllbarkeitsproblem der Aussagenlogik (engl. Satisfiability)}
    \acro{UIP}{unique implication point}
    \acro{VSIDS}{variable state independent decaying sum}
\end{acronym}


\newpage
\else
\fi
\ifcase\myIllustrationDirectory

\or
    \listoffigures
\newpage

\else
\fi
\ifcase\myTableDirectory

\or
    \listoftables
\newpage

\else
\fi
\ifcase\myListingDirectory

\or
    \section*{\langListingsverzeichnis}

\markboth{\langListingsverzeichnis}{\langListingsverzeichnis}

\addcontentsline{toc}{section}{\langListingsverzeichnis}

% Ein gemeinsames Verzeichnis für Quellcode (lstlisting) und Pseudocode (algorithm).
% Beide schreiben in die Datei .lol (siehe \ext@algorithm in vorlage.tex),
% damit die Einträge in Dokumentreihenfolge erscheinen.
\makeatletter
\@starttoc{lol}
\makeatother

\newpage

\else
\fi
\ifcase\myFormularDirectory

\or
    \newpage

\section*{\langFormelverzeichnis}

\markboth{\langFormelverzeichnis}{\langFormelverzeichnis} 

\addcontentsline{toc}{section}{\langFormelverzeichnis}

% Prefix für die Gleichung (z.B. FV. oder LF.)
\renewcommand{\theequation}{\langFormelPrefix\arabic{equation}} 
\counterwithin{equation}{section}

\newcommand{\dotfillref}[2]{#1 \dotfill #2}



\else
\fi
\ifcase\mySymbolDirectory

\or   
	\renewcommand{\nomname}{\langSymbolverzeichnis}




\printnomenclature
\else
\fi


	% Kapitel
	\pagenumbering{arabic}
	\section{Einleitung}\label{sec:einleitung}
Sudoku ist weltweit eines der bekanntesten Logikrätsel und für viele ein einfacher Zeitvertreib, stellt aber aus der Perspektive der theoretischen Informatik ein schwieriges Entscheidungsproblem dar. Bei der Lösung muss aus mehreren möglichen Belegungen eine gültige Lösung gefunden werden. Dadurch eignet sich Sudoku um auf zentrale Konzepte wie Entscheidungsprobleme, Lösbarkeit und Reduktionen einzugehen. 

\subsection{Motivation und Problemstellung}\label{sec:motivation}

Für die algorithmische Lösung von Sudokus gibt es verschiedene Herangehensweisen. Naheliegend hierbei ist das Backtracking, bei dem Felder schrittweise belegt und bei Widerspruch wieder zurückgenommen werden. Ein weiterer Ansatz ergibt sich durch das Übersetzen des Rätsels in eine aussagenlogische Formel, um diese mithilfe eines Solvers für das Erfüllbarkeitsproblem der Aussagenlogik (engl. Satisfiability, \acs{SAT})\acused{SAT} zu lösen. Diese Übersetzung ist ein Beispiel für Reduktion, bei der ein Problem in ein anderes, besser untersuchtes Problem zurückgeführt wird, um dort bereits entdeckte Lösungsverfahren nutzen zu können. In beiden Ansätzen lässt sich das Lösen von Sudokus als Entscheidungsproblem untersuchen, da geprüft wird, ob ein gegebenes Sudoku lösbar ist.
Daraus leitet sich die Frage ab, inwiefern sich die beiden Lösungsansätze unterscheiden, sowohl in Hinsicht auf ihr Vorgehen als auch in Bezug auf den notwendigen Rechenaufwand, insbesondere beim Lösen schwerer Sudokus.

\subsection{Zielsetzung}\label{sec:zielsetzung}

Ziel dieser Arbeit ist es, die beiden Lösungsansätze anhand ihrer Lösungswege, ihrer Laufzeit und der Anzahl benötigter Schritte zu vergleichen, um die Leitfrage aus Abschnitt~\ref{sec:motivation} umfangreich zu beantworten. Dazu wird das Lösen von Sudokus zunächst formal als Entscheidungsproblem beschrieben. Anschließend werden die beiden Verfahren in Python implementiert und ihre Lösungswege Schritt für Schritt in einer interaktiven Visualisierung sichtbar gemacht. Zum Schluss werden die beiden Verfahren an leichten und schweren Sudokus verglichen.


\subsection{Aufbau der Arbeit}\label{sec:aufbau}

Der Ablauf dieser Arbeit beginnt mit dem Kapitel theoretische Grundlagen, in dem auf Sudokus als Entscheidungsproblem, Komplexität und Reduktionen, Backtracking und SAT eingegangen wird, um das Fundament für die darauffolgende Arbeit zu schaffen. Daraufhin folgt das Kapitel Umsetzung in dem zunächst der Aufbau des Projekts, die beiden Lösungsverfahren, die Reduktion auf SAT und zum Schluss die Visualisierung aufgegriffen werden. 
Das Kapitel Evaluation geht auf den Versuchsaufbau und die Ergebnisse der Messungen beider Verfahren ein und führt zum Schluss mit einer Diskussion in das Fazit über. Im Fazit werden die zentralen Ergebnisse dieser Arbeit zusammengefasst, die Leitfrage beantwortet und abschließend ein Ausblick gegeben.


\section{Theoretische Grundlagen}\label{sec:grundlagen}

\subsection{Sudoku als Entscheidungsproblem}\label{sec:csp}
Ein klassisches Sudoku besteht aus einem Gitter mit $9 \times 9$ Zellen, das in neun Zeilen, neun Spalten und neun Blöcke der Größe $3 \times 3$ unterteilt ist. Einige Zellen sind bereits zu Beginn mit Ziffern belegt und werden als Vorgaben bezeichnet. Das Sudoku gilt als gelöst, wenn jede Zelle eine Ziffer von 1 bis 9 enthält und jede Ziffer pro Zeile, Spalte und Block genau einmal vorkommt.

Aus Sicht der theoretischen Informatik lässt sich diese Aufgabe auf zwei Arten formulieren. Als Suchproblem wird eine konkrete Lösung verlangt. 
Als Entscheidungsproblem wird lediglich gefragt, ob zu einem vorgegebenen, teilweise gefüllten Sudoku eine gültige Lösung existiert, die Antwort lautet also entweder Ja oder Nein. Die Komplexitätstheorie betrachtet in der Regel Entscheidungsprobleme, weshalb diese Formulierung die Grundlage für Abschnitt~\ref{sec:komplexitaet} bildet. Beide Sichtweisen hängen eng zusammen, da die in dieser Arbeit betrachteten Verfahren die Frage nach der Lösbarkeit beantworten, indem sie versuchen eine Lösung zu konstruieren.

Formal werden Entscheidungsprobleme als Sprachen beschrieben. Eine Sprache ist dabei eine Menge von Wörtern, die aus den Zeichen eines festgelegten Alphabets gebildet werden \vgl{sipser2013}. Ein Alphabet bezeichnet eine endliche Menge von benötigten Zeichen, bei einem Sudoku sind die Zeichen die Ziffern von 0 bis 9. Ein Sudoku lässt sich als Wort der Länge 81 darstellen, indem die Zellen zeilenweise aneinandergereiht werden und die Ziffer 0 eine leere Zelle kennzeichnet. Die Menge aller lösbaren Sudokus bildet die Sprache.

Die Frage, ob ein gegebenes Sudoku lösbar ist, entspricht damit dem Wortproblem, also der Frage, ob das Wort zu der Sprache gehört.

Ein übliches $9 \times 9$ Sudoku kann genau eine Lösung haben, wenn 17 oder mehr Vorgaben gegeben sind \vgl{mcguire2014}. Sudokus mit wenigen oder ungünstig verteilten Vorgaben können dagegen mehrere Lösungen zulassen, widersprüchliche Vorgaben keine einzige. Für das Entscheidungsproblem ist nur relevant, ob mindestens eine zulässige Lösung existiert.

\subsection{Komplexität, Lösbarkeit und Reduktionen}\label{sec:komplexitaet}
Nicht jedes Entscheidungsproblem lässt sich algorithmisch beantworten. Turing zeigte mithilfe des Halteproblems, dass es Fragen gibt, für die kein Algorithmus existiert, der für jede Eingabe nach endlich vielen Schritten die richtige Antwort liefert \vgl{sipser2013}. Solche Probleme heißen unentscheidbar. Das Sudoku-Problem aus Abschnitt~\ref{sec:csp} gehört nicht dazu, da der Suchraum eines Sudokus endlich ist. Bei $k$ leeren Zellen existieren höchstens $9^k$ Belegungen, die sich der Reihe nach überprüfen lassen. Ein solches Verfahren findet entweder eine Lösung oder stellt fest, dass hierbei keine Lösung existiert. Sudoku ist somit also entscheidbar.

Die Komplexitätstheorie ordnet Entscheidungsprobleme nach diesem Aufwand in Klassen ein, gemessen in Abhängigkeit von der Problemgröße $n$. Die Klasse P enthält alle Probleme, die ein deterministischer Algorithmus in polynomieller Zeit löst. Die Klasse NP enthält alle Probleme, bei denen sich eine vorgeschlagene Lösung in polynomieller Zeit überprüfen lässt \vgl{sipser2013}. Gleichwertig lässt sich NP über nichtdeterministische Algorithmen beschreiben, die eine Lösung raten und anschließend in polynomieller Zeit prüfen. Jedes Problem aus P liegt auch in NP. Sudoku liegt in NP, da für ein vollständig gefülltes Sudoku alle Zeilen, Spalten und Blöcke auf doppelte Ziffern nur geprüft werden müssen.

Um Probleme hinsichtlich ihrer Schwierigkeit zu vergleichen, dient die polynomielle Reduktion. Ein Problem $A$ ist auf ein Problem $B$ reduzierbar, geschrieben $A \leq_p B$, wenn sich jede Eingabe von $A$ in polynomieller Zeit so in eine Eingabe von $B$ umformen lässt, dass beide dieselbe Antwort besitzen \vgl{sipser2013}. Ein Verfahren für $B$ löst damit auch $A$. Ein Problem heißt NP-vollständig, wenn es in NP liegt und sich jedes Problem aus NP darauf reduzieren lässt. Cook wies als erstes Problem das Erfüllbarkeitsproblem der Aussagenlogik als NP-vollständig nach \vgl{cook1971}. Für NP-vollständige Probleme ist kein Algorithmus mit polynomieller Laufzeit bekannt, ob ein solcher existiert, ist als P-NP-Problem eine der bekanntesten offenen Fragen der Informatik und eines der Millenial-Probleme. Die bekannten exakten Verfahren benötigen im ungünstigsten Fall exponentielle Laufzeit, weshalb solche Probleme bei einer großen Menge an Eingaben als praktisch nicht handhabbar gelten.

Für Sudoku gilt diese Einordnung in verallgemeinerter Form. Auf Sudokus der Größe $n^2 \times n^2$ ist die Frage nach der Lösbarkeit NP-vollständig \vgl{yato2003}. Das klassische Sudoku mit $n = 3$ besitzt eine feste Größe und ist damit streng genommen in konstanter Zeit lösbar. Die Komplexität des allgemeinen Problems erklärt dennoch, warum der Aufwand einfacher Suchverfahren bei ungünstigen Rätseln stark ansteigt.

Da Sudoku in NP liegt und SAT NP-vollständig ist, existiert eine polynomielle Reduktion von Sudoku auf SAT. Für ein Sudoku mit Kantenlänge $m$ genügen dabei $m^3$ Variablen und $O(m^4)$ Klauseln. Diese Reduktion bildet die Grundlage des zweiten Lösungsansatzes dieser Arbeit und wird in Abschnitt~\ref{sec:impl-sat} konkret umgesetzt.



\subsection{Backtracking und Heuristiken}\label{sec:backtracking}


\subsection{Aussagenlogik und das SAT-Problem}\label{sec:sat}


\subsection{Der DPLL-Algorithmus}\label{sec:dpll}

% Code:
%   src/sudoku_solver/solvers/backtracking/
%     naive.py       (solve, stack_solve, solve_steps)
%     heuristic.py   (MRV, Forward Checking)
%     counting.py    (Lösungsanzahl, Eindeutigkeit, Baumgröße)
%     estimation.py  (Knuth-Schätzer)
%   Zahlen für die Tabellen:
%   python -m sudoku_solver.solvers.backtracking.report


\section{Backtracking-Solver}\label{sec:bt}

Als erster der beiden Lösungsansätze wird in diesem Kapitel der Backtracker behandelt.
Hierfür wird das Sudoku als Entscheidungsproblem betrachtet, jedoch ist Backtracking ein
Suchverfahren, welches Schritt für Schritt eine konkrete Lösung konstruiert. 
Dafür wird das Sudoku als ein Suchproblem betrachtet, dessen Teillösungen einen Suchbaum bilden.

Als erstes werden die Grundlagen für das Entscheidungsproblem gelegt (Abschnitt \ref{sec:bt-suchbaum}),
anschließend wird der Algorithmus und seine Implementierung beschrieben (Abschnitt \ref{sec:bt-algorithmus}).
Daraufhin wird begründet, dass das Verfahren korrekt, vollständig und stets terminiert ist (Abschnitt \ref{sec:bt-korrektheit}),
und sein Bedarf an Laufzeit und Speicher eingeordnet (Abschnitt \ref{sec:bt-komplexitaet}).
Die letzten beiden Abschnitte behandeln Heuristiken um den Suchbaum zu verkleinern (Abschnitt \ref{sec:bt-heuristiken}),
sowie abschließend ein Verfahren nach Knuth, das die Größe des Suchbaums vor dem durchsuchen schätzt (Abschnitt \ref{sec:bt-knuth}). 

% ----------------------------------------------------------------------------
\subsection{Zustandsraum und Suchbaum}\label{sec:bt-suchbaum}

Grundlage des Backtrackings ist ein Zustandsraum, der die schrittweise entstehenden Teillösungen eines Sudokus beschreibt.
Dazu werden die $k$ leeren Zellen zeilenweise von links oben nach rechts unten als $z_1, \dots, z_k$ nummeriert. 
Ein \emph{Zustand} ist eine partielle Belegung, die den ersten $i$ dieser Zellen mit $0 \leq i \leq k$ 
jeweils eine Ziffer aus $\{1, \dots, 9\}$ zuordnet. 
Er heißt \emph{konsistent}, wenn zusammen mit den Vorgaben keine Zeile, keine Spalte und kein Block eine Ziffer mehrfach enthält. 
Eine Lösung ist damit ein konsistenter Zustand, der alle $k$ leeren Zellen belegt.

Die konsistenten Zustände bilden einen \emph{Suchbaum}. 
Seine Wurzel ist das Sudoku mit seinen Vorgaben, 
und die Kinder eines Knotens der Tiefe $i$ sind alle konsistenten Zustände, die ihn um eine Ziffer für $z_{i+1}$ erweitern. 
Jeder Knoten außer der Wurzel entspricht somit genau einer gesetzten Ziffer. 
Der Verzweigungsgrad, also die Anzahl der zulässigen Ziffern der nächsten Zelle, beträgt höchstens neun, 
die Tiefe des Baums höchstens $k$. 
Ein Blatt ist entweder eine Lösung oder eine \emph{Sackgasse}, 
also ein konsistenter Zustand, für dessen nächste Zelle keine Ziffer mehr zulässig ist.

Wie in Abschnitt~\ref{sec:komplexitaet} beschrieben, 
ließe sich die Lösbarkeit auch durch Prüfen aller $9^k$ vollständigen Belegungen entscheiden. 
Bei $k = 60$ leeren Zellen wären das jedoch etwa $1{,}8 \cdot 10^{57}$ Belegungen, 
für die selbst eine Milliarde Prüfungen pro Sekunde mehr als $10^{40}$ Jahre bräuchten. 
Backtracking vermeidet diesen Aufwand, indem es die Konsistenz bereits nach jeder gesetzten Ziffer prüft \vgl{russell2021}. 
Grundlage ist eine Monotonie-Eigenschaft: Enthält ein Zustand einen Widerspruch, so enthält ihn auch jede seiner Erweiterungen. 
Ein inkonsistenter Zustand wird daher gar nicht erst erzeugt, 
und mit ihm entfällt sein gesamter Teilbaum mit bis zu $9^{k-i}$ Belegungen. 
Dieses Abschneiden wird als \emph{Pruning} bezeichnet.

Backtracking ist also eine Tiefensuche auf dem Suchbaum; 
Der aktuelle Zustand wird so lange um zulässige Ziffern erweitert, bis eine Lösung oder eine Sackgasse erreicht ist. 
In einer Sackgasse wird die zuletzt gesetzte Ziffer zurückgenommen und im Elternknoten die nächste Alternative probiert. 
Diesem Rücksprung verdankt das Verfahren seinen Namen ("Backtracking"). 
Der Suchbaum wird dabei nie vollständig gespeichert, sondern während der Suche schrittweise erzeugt. 
Abbildung~\ref{fig:bt-suchbaum} zeigt dies am vollständigen Suchbaum eines $4 \times 4$-Sudokus 
mit den Ziffern 1 bis 4 und $k = 5$ leeren Zellen: 
Statt der $4^5 = 1024$ möglichen Belegungen setzt das Backtracking nur sieben Ziffern 
und durchläuft zwei Sackgassen, bevor es die Lösung findet.

\begin{figure}[hbt]
	\centering
	\begin{minipage}[c]{0.27\textwidth}
		\centering
		\begin{tikzpicture}[scale=0.6]
			\draw[gray] (0,0) grid (4,4);
			\draw[very thick] (0,0) rectangle (4,4);
			\draw[very thick] (2,0) -- (2,4);
			\draw[very thick] (0,2) -- (4,2);
			% Vorgaben, Zeile für Zeile von oben
			\foreach \c/\z in {0/1, 1/2, 2/3, 3/4} \node at (\c+0.5, 3.5) {\textbf{\z}};
			\foreach \c/\z in {0/3, 1/4} \node at (\c+0.5, 2.5) {\textbf{\z}};
			\foreach \c/\z in {0/4, 3/2} \node at (\c+0.5, 1.5) {\textbf{\z}};
			\foreach \c/\z in {0/2, 2/4, 3/3} \node at (\c+0.5, 0.5) {\textbf{\z}};
			% leere Zellen
			\node[gray] at (2.5, 2.5) {\footnotesize $z_1$};
			\node[gray] at (3.5, 2.5) {\footnotesize $z_2$};
			\node[gray] at (1.5, 1.5) {\footnotesize $z_3$};
			\node[gray] at (2.5, 1.5) {\footnotesize $z_4$};
			\node[gray] at (1.5, 0.5) {\footnotesize $z_5$};
		\end{tikzpicture}
	\end{minipage}%
	\hfill
	\begin{minipage}[c]{0.7\textwidth}
		\centering
		\begin{forest}
			for tree={
				draw, rounded corners=2pt, font=\footnotesize, inner sep=2pt,
				l sep=5mm, s sep=8mm, edge={->}
			}
			[Vorgaben, name=wurzel, label={[font=\tiny, text=gray]right:0}
				[{$z_1 = 1$}, name=sack1, fill=red!15, label={[font=\tiny, text=gray]right:1}]
				[{$z_1 = 2$}, label={[font=\tiny, text=gray]right:2}
					[{$z_2 = 1$}, name=knoten, label={[font=\tiny, text=gray]right:3}
						[{$z_3 = 1$}, name=sack2, fill=red!15, label={[font=\tiny, text=gray]right:4}]
						[{$z_3 = 3$}, label={[font=\tiny, text=gray]right:5}
							[{$z_4 = 1$}, label={[font=\tiny, text=gray]right:6}
								[{$z_5 = 1$}, fill=green!20, label={[font=\tiny, text=gray]right:7}]
							]
						]
					]
				]
			]
			\draw[->, dashed, red!70!black] (sack1.north west) to[bend left=30] (wurzel.west);
			\draw[->, dashed, red!70!black] (sack2.north west) to[bend left=30] (knoten.west);
		\end{forest}
	\end{minipage}
	\caption[Suchbaum des Backtrackings für ein $4 \times 4$-Sudoku]{Links ein $4 \times 4$-Sudoku mit Vorgaben (fett) und den leeren Zellen $z_1$ bis $z_5$, 
	rechts sein vollständiger Suchbaum. Rot markiert sind Sackgassen, grün die Lösung, 
	gestrichelt die Rücksprünge. Die grauen Zahlen geben die Reihenfolge an, 
	in der die Tiefensuche die Knoten besucht.}
	\label{fig:bt-suchbaum}
\end{figure}

% ----------------------------------------------------------------------------
\subsection{Algorithmus und Implementierung}\label{sec:bt-algorithmus}

Die Implementierung im Modul \texttt{naive.py} umfasst eine rekursive und eine iterative Variante 
desselben Verfahrens sowie einen Generator für die Visualisierung. 
Alle durchlaufen die Zellen zeilenweise, probieren die Ziffern aufsteigend von 1 bis 9 
und durchsuchen damit denselben Suchbaum in derselben Reihenfolge. 
Zur Abgrenzung von den Heuristiken in Abschnitt~\ref{sec:bt-heuristiken} 
wird diese Vorgehensweise als \emph{naives Backtracking} bezeichnet.

\begin{algorithm}[hbt]
	\caption{Rekursives Backtracking (\texttt{solve})}\label{alg:bt-rekursiv}
	\begin{algorithmic}[1]
		\Function{Solve}{$G, r, c$} \Comment{Gitter $G$, aktuelle Zelle $(r, c)$, Indizes 0 bis 8}
			\If{$r = 9$} \Return \textbf{true} \Comment{alle Zeilen belegt}
			\EndIf
			\If{$c = 9$} \Return \Call{Solve}{$G, r + 1, 0$} \Comment{weiter in der nächsten Zeile}
			\EndIf
			\If{$G[r][c] \neq 0$} \Return \Call{Solve}{$G, r, c + 1$} \Comment{Vorgabe überspringen}
			\EndIf
			\For{$z \gets 1, \dots, 9$}
				\If{\Call{Valid}{$G, r, c, z$}}
					\State $G[r][c] \gets z$
					\If{\Call{Solve}{$G, r, c + 1$}} \Return \textbf{true}
					\EndIf
					\State $G[r][c] \gets 0$ \Comment{Rücksprung}
				\EndIf
			\EndFor
			\State \Return \textbf{false}
		\EndFunction
	\end{algorithmic}
\end{algorithm}

Die rekursive Variante (Algorithmus~\ref{alg:bt-rekursiv}) erhält das Gitter und die aktuelle Zelle. 
Sie überspringt Vorgaben, wechselt am Zeilenende in die nächste Zeile und meldet Erfolg, 
sobald sie hinter der letzten Zeile angekommen ist. 
In einer leeren Zelle probiert sie nacheinander jede zulässige Ziffer und setzt die Suche rekursiv mit der nächsten Zelle fort. 
Scheitert dieser Aufruf, wird die Zelle wieder geleert (Zeile~9). 
Sind alle Ziffern erfolglos, gibt der Aufruf das Scheitern an seinen Aufrufer zurück, 
der daraufhin seine nächste Ziffer probiert. 
Die Zulässigkeitsprüfung \textsc{Valid} (\texttt{is\_valid\_number}) betrachtet nur Zeile, Spalte und Block der neuen Zelle: 
Da der bisherige Zustand konsistent ist, kann ein Widerspruch nur dort entstehen. 
Der Prüfaufwand pro Knoten ist damit konstant. 
Vor der Suche wird mit \texttt{is\_valid} einmalig geprüft, ob die Vorgaben widerspruchsfrei sind 
(siehe Abschnitt~\ref{sec:bt-korrektheit}).

Bei der Rekursion liegt der Zustand der Suche implizit auf dem Call-Stack, 
wobei jeder aktive Aufruf einem Knoten auf dem Pfad von der Wurzel zum aktuellen Zustand entspricht. 
Die iterative Variante \texttt{stack\_solve} macht diesen Stack explizit und ersetzt die Rekursion durch eine Schleife. 
Dazu werden die leeren Zellen $z_1, \dots, z_k$ vorab bestimmt. 
Ein Stackeintrag $(i, z)$ hält fest, dass die Zelle $z_i$ zuletzt die Ziffer $z$ erhalten hat. 
In jedem Durchlauf wird der oberste Eintrag entnommen und ab der Ziffer $z + 1$ nach der nächsten zulässigen Ziffer gesucht. 
Gelingt dies, kommt der aktualisierte Eintrag zurück auf den Stack und darüber ein neuer Eintrag für die Zelle $z_{i+1}$. 
Andernfalls bleibt der Eintrag entfernt, und der nächste Durchlauf setzt bei der vorherigen Zelle fort, 
was dem Rücksprung entspricht. 
Ist der Stack leer, sind alle Alternativen erschöpft und das Sudoku ist unlösbar.

Der eigen implementierte, explizite Stack hat die zwei folgenden Vorteile: 
Zum einen ist die Suche nicht durch die maximale Rekursionstiefe begrenzt, was bei $9 \times 9$-Sudokus zwar unkritisch ist, 
bei größeren Varianten aber relevant wird. 
Vor allem aber liegt der Suchzustand jederzeit vollständig vor, sodass die Suche nach jedem Schritt unterbrochen werden kann. 
Der Generator \texttt{solve\_steps} baut auf dieses Prinzip auf. 
Die Funktion gibt für jeden Schritt die jeweilige Aktion zurück mitsamt einer Kopie des Gitters. 

% ----------------------------------------------------------------------------
\subsection{Korrektheit, Vollständigkeit und Terminierung}\label{sec:bt-korrektheit}

Um die Lösbarkeit zu entscheiden, muss der Backtracker korrekt, vollständig und terminierend sein.

\emph{Korrektheit:} Da nur Ziffern gesetzt werden, die \textsc{Valid} erlaubt, 
und die Vorgaben vorab mit \texttt{is\_valid} geprüft werden, ist jeder erreichte Zustand konsistent. 
Meldet der Algorithmus Erfolg, sind alle $k$ Zellen belegt, es liegt also eine Lösung vor. 
Ohne die Vorabprüfung würde ein vollständig ausgefülltes, aber ungültiges Gitter als lösbar gemeldet.

\emph{Vollständigkeit:} Wegen der Monotonie aus Abschnitt~\ref{sec:bt-suchbaum} 
verwirft das Pruning nur Teilbäume ohne Lösung. Alle übrigen Knoten durchläuft die Tiefensuche systematisch. 
Eine existierende Lösung wird daher gefunden, 
und dass das Sudoku andernfalls unlösbar ist wird erst nach Durchsuchen des gesamten Baums gemeldet.

\emph{Terminierung:} Der Suchbaum ist endlich, da seine Tiefe höchstens $k$ und sein Verzweigungsgrad höchstens neun beträgt. 
Weil jede Zelle ihre Ziffern aufsteigend ab der zuletzt probierten durchläuft, wird kein Zustand doppelt erzeugt.

Der Backtracker hält somit auf jeder Eingabe mit der korrekten Antwort 
und ist ein Entscheidungsverfahren für die Sprache der lösbaren Sudokus aus Abschnitt~\ref{sec:csp}.

% ----------------------------------------------------------------------------
\subsection{Laufzeit- und Speicherkomplexität}\label{sec:bt-komplexitaet}

Da der Prüfaufwand pro Knoten konstant ist (Abschnitt~\ref{sec:bt-algorithmus}), 
bestimmt die Anzahl der besuchten Knoten die Laufzeit. 
Im ungünstigsten Fall wird der gesamte Suchbaum mit höchstens
\begin{equation}
	\sum_{i=0}^{k} 9^i = \frac{9^{k+1} - 1}{8} \in O(9^k)
\end{equation}
Knoten durchlaufen. Für das $9 \times 9$-Sudoku fester Größe ist diese Aussage 
nur eingeschränkt aussagekräftig (Abschnitt~\ref{sec:komplexitaet}). 
Beim verallgemeinerten $n^2 \times n^2$-Sudoku beträgt der Verzweigungsgrad bis zu $n^2$, 
die Laufzeit liegt also in $O(n^{2k})$ und wächst exponentiell mit der Anzahl der leeren Zellen. 
Wie viele Knoten tatsächlich besucht werden, hängt von der Instanz und der Suchreihenfolge ab 
und wird in Kapitel~\ref{sec:evaluation} empirisch untersucht.

Der Speicherbedarf ist dagegen gering: Neben dem Gitter wird nur der Pfad von der Wurzel zum aktuellen Knoten 
mit höchstens $k$ Stackeinträgen gespeichert, also $O(k)$. 
Eine Breitensuche müsste dagegen ganze Ebenen mit bis zu $9^i$ Zuständen vorhalten.

Dem exponentiellen Suchaufwand steht die Überprüfung einer Lösung gegenüber, 
die \texttt{is\_solved} in polynomieller Zeit leistet. 
Diese Asymmetrie entspricht der Einordnung in NP aus Abschnitt~\ref{sec:komplexitaet}: 
Das Backtracking ersetzt das nichtdeterministische Raten einer Lösung durch systematisches Ausprobieren 
und bezahlt dies mit exponentiellem Aufwand. 
Da das verallgemeinerte Sudoku NP-vollständig ist \vgl{yato2003}, 
existiert unter der Annahme $\mathrm{P} \neq \mathrm{NP}$ kein Algorithmus mit polynomieller Laufzeit. 
Der exponentielle Worst Case ist somit eine Eigenschaft des Problems und nicht des Verfahrens. 
Verbessern lässt sich jedoch der Aufwand typischer Instanzen durch eine geschicktere Suchreihenfolge, 
wie der folgende Abschnitt zeigt.

% ----------------------------------------------------------------------------
\subsection{Heuristiken zur Verkleinerung des Suchbaums}\label{sec:bt-heuristiken}

Am exponentiellen Worst Case lässt sich nach Abschnitt~\ref{sec:bt-komplexitaet} nichts ändern. 
Der Aufwand für eine konkrete Instanz hängt jedoch stark davon ab, in welcher Reihenfolge die Zellen belegt werden. 
Das naive Backtracking belegt sie in einer festen, zeilenweisen Reihenfolge, unabhängig vom aktuellen Zustand. 
Im Modul \texttt{heuristic.py} sind zwei Heuristiken aus dem Bereich der Constraint-Satisfaction-Probleme umgesetzt, 
die die Suche an den aktuellen Zustand anpassen \vgl{russell2021}. 
Beide lassen sich unabhängig voneinander zuschalten.

\emph{\ac{MRV}:} Statt der nächsten Zelle in fester Reihenfolge wird die leere Zelle 
mit den wenigsten zulässigen Ziffern, den sogenannten Kandidaten, als Nächstes belegt. 
Dahinter steht das Fail-First-Prinzip \vgl{haralick1980}: 
Muss ein Zweig scheitern, soll er möglichst früh und damit möglichst nah an der Wurzel scheitern, 
da so der größte Teilbaum eingespart wird. 
Hat eine Zelle nur einen Kandidaten, ist ihre Belegung erzwungen und der Baum verzweigt sich an dieser Stelle nicht. 
Hat eine Zelle keinen Kandidaten, wird die Sackgasse sofort erkannt, statt zuvor noch weitere Zellen zu belegen. 
Der Suchbaum ist weiterhin wie in Abschnitt~\ref{sec:bt-suchbaum} definiert, 
nur wird die nächste Zelle $z_{i+1}$ nun abhängig vom aktuellen Zustand gewählt. 
Korrektheit, Vollständigkeit und Terminierung aus Abschnitt~\ref{sec:bt-korrektheit} bleiben daher erhalten.

\emph{\ac{FC}:} Nach jeder gesetzten Ziffer wird geprüft, 
ob in der Zeile, der Spalte oder dem Block der Zelle eine leere Zelle ohne Kandidaten entstanden ist. 
In diesem Fall wird die Ziffer sofort zurückgenommen, 
auch wenn die betroffene Zelle in der festen Reihenfolge erst viel später an der Reihe wäre.

\begin{table}[hbt]
	\centering
	\caption[Gesetzte Ziffern bis zur ersten Lösung: naives Backtracking und Heuristiken]{Gesetzte Ziffern bis zur ersten Lösung für das naive Backtracking und die Heuristiken.
	Die Suche wurde nach $500\,000$ gesetzten Ziffern abgebrochen.} {\emph{Eigene Messung}}
	\label{tab:bt-heuristiken}
	\begin{tabular}{lrrrrr}
		\toprule
		Sudoku & leere Zellen & naiv & naiv + \acs{FC} & \acs{MRV} & \acs{MRV} + \acs{FC} \\
		\midrule
		Wikipedia                      & 51 & $4\,208$     & $516$        & $51$       & $51$ \\
		AI Escargot                    & 58 & $8\,969$     & $5\,421$     & $217$      & $217$ \\
		Inkala 2012                    & 60 & $49\,558$    & $22\,067$    & $10\,101$  & $10\,101$ \\
		Norvig \glqq hard1\grqq{}      & 64 & $>500\,000$  & $>500\,000$  & $481$      & $481$ \\
		Anti-Backtracking              & 64 & $>500\,000$  & $>500\,000$  & $45\,267$  & $45\,267$ \\
		Leicht für Backtracking	       & 57 & $992$        & $336$        & $3\,220$   & $3\,220$ \\
		Unlösbares Sudoku              & 72 & $>500\,000$  & $>500\,000$  & $0$        & $0$ \\
		\bottomrule
	\end{tabular}
\end{table}

Tabelle~\ref{tab:bt-heuristiken} vergleicht die Anzahl gesetzter Ziffern, also der besuchten Knoten, 
für einige Sudokus aus der Sammlung des Projekts. 
Die Werte lassen sich mit dem Skript \texttt{report.py} reproduzieren. 
Beim Wikipedia-Sudoku setzt \ac{MRV} genau 51 Ziffern für 51 leere Zellen: 
Jede Belegung ist erzwungen, und der Suchbaum besitzt keine einzige Verzweigung. 
Beim Sudoku \glqq hard1\grqq{} sinkt der Aufwand von rund 9,7 Millionen gesetzten Ziffern 
ohne Abbruch auf 481, beim Anti-Backtracking-Sudoku von rund 69 Millionen auf etwa $45\,000$. 
Am deutlichsten zeigt sich das Fail-First-Prinzip beim unlösbaren Sudoku. 
Seine letzte Zelle besitzt von Beginn an keinen Kandidaten. 
\ac{MRV} erkennt das, ohne eine einzige Ziffer zu setzen, 
während das naive Verfahren diese Zelle erst erreicht, wenn alle 71 davorliegenden Zellen belegt sind, und dies für jede mögliche Belegung erneut.

Zwei Beobachtungen fallen auf. 
Erstens setzt \ac{MRV} mit \ac{FC} in allen Fällen genau so viele Ziffern wie \ac{MRV} allein. 
Eine Zelle ohne Kandidaten, die \ac{FC} erkennen würde, ist genau die Zelle, die \ac{MRV} als Nächstes wählt 
und an der die Suche ohne weitere Platzierung scheitert. 
\ac{FC} lohnt sich daher nur bei einer festen Reihenfolge. 
Zweitens ist eine Heuristik keine Garantie: Beim Sudoku \glqq Leicht für Backtracking\grqq{} 
benötigt \ac{MRV} mit $3\,220$ gesetzten Ziffern mehr als das Dreifache des naiven Verfahrens. 
\ac{MRV} entscheidet nur nach einem lokalen Kriterium und kann dadurch in größere Teilbäume geraten, 
die die feste Reihenfolge zufällig umgeht. 
Zudem ist ein Knoten bei \ac{MRV} teurer, da die Kandidaten aller leeren Zellen bestimmt werden müssen. 
Wie sich das auf die Laufzeit auswirkt, wird in Kapitel~\ref{sec:evaluation} untersucht.

% TODO: ref auf das Laufzeit Kapitel wenn es soweit ist
Aus Sicht der theoretischen Informatik ändern die Heuristiken nur die Reihenfolge der Suche, nicht den Suchraum. 
Die obere Schranke $O(9^k)$ aus Abschnitt~\ref{sec:bt-komplexitaet} bleibt bestehen, 
und unter der Annahme $\mathrm{P} \neq \mathrm{NP}$ kann keine Heuristik die Laufzeit allgemein polynomiell machen. 
Heuristiken verkleinern den Suchbaum typischer Instanzen um Größenordnungen, die Komplexitätsklasse des Problems ändern sie nicht.

% ----------------------------------------------------------------------------
\subsection{Abschätzung der Suchbaumgröße nach Knuth}\label{sec:bt-knuth}

Tabelle~\ref{tab:bt-heuristiken} zeigt, dass sich der Aufwand selbst bei ähnlich vielen leeren Zellen 
um mehrere Größenordnungen unterscheiden kann. 
Ob eine Suche nach Millisekunden oder erst nach Jahren endet, ist vorab meist nicht erkennbar. 
\textcite{knuth1975} schlug dafür ein Monte-Carlo-Verfahren vor, das die Größe eines Backtracking-Suchbaums schätzt, 
ohne ihn zu durchlaufen. Es ist im Modul \texttt{estimation.py} umgesetzt.

Das Verfahren folgt einem zufälligen Pfad von der Wurzel bis zu einem Blatt 
und wählt in jedem Knoten eines der Kinder gleichverteilt aus. 
Bezeichnet $d_j$ die Anzahl der Kinder des $j$-ten Knotens auf diesem Pfad, 
beginnend mit der Wurzel, so ist die Schätzung für einen Pfad mit $t$ inneren Knoten
\begin{equation}
	X = 1 + d_1 + d_1 d_2 + \dots + d_1 d_2 \cdots d_t.
\end{equation}
Die Intuition dahinter ist, dass alle Geschwister eines besuchten Knotens als gleich groß angenommen werden. 
Hat die Wurzel zwei Kinder, werden auf der ersten Ebene zwei Knoten geschätzt. 
Hat das gewählte Kind wieder zwei Kinder, werden auf der zweiten Ebene $2 \cdot 2 = 4$ Knoten geschätzt. 
Als Gesamtschätzung dient der Mittelwert über viele unabhängige Pfade.

Diese Schätzung ist erwartungstreu. 
Ein Knoten $v$ in Tiefe $j$ liegt mit der Wahrscheinlichkeit $p_v = 1/(d_1 \cdots d_j)$ auf dem zufälligen Pfad 
und trägt in diesem Fall genau den Summanden $d_1 \cdots d_j = 1/p_v$ zu $X$ bei. 
Für den Erwartungswert über alle Knoten $v$ des Suchbaums $T$ folgt daher
\begin{equation}
	E[X] = \sum_{v \in T} p_v \cdot \frac{1}{p_v} = |T|.
\end{equation}
Für den Suchbaum aus Abbildung~\ref{fig:bt-suchbaum} gibt es folgende drei mögliche Pfade. 
Der Pfad über $z_1 = 1$ wird mit Wahrscheinlichkeit $\frac{1}{2}$ gewählt und liefert $X = 1 + 2 = 3$, 
die beiden Pfade über $z_3 = 1$ und $z_3 = 3$ mit Wahrscheinlichkeit je $\frac{1}{4}$ 
liefern $X = 9$ und $X = 17$. 
Der Erwartungswert $\frac{1}{2} \cdot 3 + \frac{1}{4} \cdot 9 + \frac{1}{4} \cdot 17 = 8$ 
entspricht genau den acht Knoten des Baums.

\begin{table}[hbt]
	\centering
	\caption[Exakte Suchbaumgröße und Schätzung nach Knuth]{Exakte Größe des vollständigen Suchbaums und Schätzung nach Knuth aus $1000$ Pfaden.
	Die exakte Zählung wurde nach $500\,000$ Knoten abgebrochen.} {\emph{Eigene Messung}}
	\label{tab:bt-knuth}
	\small
	\setlength{\tabcolsep}{4pt}
	\begin{tabular}{llrrrrr}
		\toprule
		Sudoku & Reihenfolge & exakt & Schätzung & Std.-Fehler & min. $X$ & max. $X$ \\
		\midrule
		Wikipedia         & naiv     & $4\,632$     & $5\,546$      & $1\,187$      & $64$      & $722\,350$ \\
		                  & \acs{MRV} & $52$        & $52$          & $0$           & $52$      & $52$ \\
		AI Escargot       & naiv     & $18\,885$    & $40\,462$     & $25\,714$     & $65$      & $2{,}5 \cdot 10^{7}$ \\
		                  & \acs{MRV} & $2\,232$    & $2\,088$      & $150$         & $33$      & $55\,433$ \\
		Anti-Backtracking & naiv     & $>500\,000$  & $5{,}2 \cdot 10^{7}$ & $2{,}3 \cdot 10^{7}$ & $8\,282$ & $1{,}6 \cdot 10^{10}$ \\
		                  & \acs{MRV} & $73\,483$   & $64\,495$     & $9\,013$      & $46$      & $4{,}3 \cdot 10^{6}$ \\
		\bottomrule
	\end{tabular}
\end{table}

Tabelle~\ref{tab:bt-knuth} vergleicht die Schätzung mit der exakten Größe des vollständigen Suchbaums, 
die durch Zählen aller Knoten mit \texttt{tree\_size} bestimmt wurde. 
Anders als in Tabelle~\ref{tab:bt-heuristiken} umfasst der vollständige Baum auch alle Knoten nach der ersten Lösung. 
Er entspricht damit dem Aufwand einer erschöpfenden Suche, wie sie bei unlösbaren Sudokus nötig ist, aber auch beim Zählen aller Lösungen oder beim Prüfen, ob eine Lösung eindeutig ist. 
Für die Bäume mit \ac{MRV} liegt die Schätzung nahe am exakten Wert. 
Beim Wikipedia-Sudoku verzweigt der Baum nie, sodass jeder Pfad den exakten Wert liefert.

In unbalancierten Bäumen, wie sie das naive Verfahren erzeugt, schwanken die einzelnen Schätzungen dagegen stark. 
Die meisten Pfade enden schnell in einer Sackgasse, 
während seltene Pfade in einen riesigen Teilbaum führen und dann extrem große Werte liefern. 
Die Verteilung von $X$ ist dadurch stark rechtsschief, 
und der Standardfehler des Mittelwerts sinkt nur mit $1/\sqrt{N}$ in der Anzahl $N$ der Pfade. 
Beim Sudoku AI Escargot überschätzt der Mittelwert die exakte Größe um mehr als das Doppelte, 
weil ein einzelner Pfad allein rund 25 Millionen Knoten schätzt. 
Beim Anti-Backtracking-Sudoku reichen die Einzelschätzungen über sechs Größenordnungen. 
Der Mittelwert von rund 52 Millionen liegt sogar unter den 69 Millionen Knoten, 
die das naive Verfahren bereits bis zur ersten Lösung besucht. 
Er unterschätzt also deutlich, weil die seltenen, sehr großen Werte bei $1000$ Pfaden kaum gezogen werden.

Damit zeigt sich die zentrale Schwäche des Verfahrens: 
Gerade bei sehr großen Bäumen, für die eine Vorhersage am wertvollsten wäre, ist die Schätzung am unzuverlässigsten. 
Für eine Einordnung der Größenordnung genügt sie dennoch. 
Für das unlösbare Sudoku schätzt das Verfahren in naiver Reihenfolge rund $8 \cdot 10^{16}$ Knoten 
und sagt damit richtig voraus, dass die naive Suche praktisch nicht endet, 
während der Baum mit \ac{MRV} aus nur einem Knoten besteht.

% ----------------------------------------------------------------------------

\section{SAT-basierter Solver} \label{sec:sat}

Dieses Kapitel beschreibt einen Sudoku-Solver, der das Rätsel als aussagenlogische Formel formuliert und mit einem \acs{SAT}-Solver löst.
Das Lösen durchläuft drei Schritte: Das Gitter wird in Klauseln kodiert, ein \acs{SAT}-Solver sucht ein Modell, und dieses Modell wird zurück in ein ausgefülltes Gitter übersetzt.
Abschnitt~\ref{sec:sat-encoding} führt \acs{SAT} ein und beschreibt die Kodierung, Abschnitt~\ref{sec:algorithms} die Suchalgorithmen \acl{DPLL} (\acs{DPLL})\acused{DPLL} und \acl{CDCL} (\acs{CDCL})\acused{CDCL} und Abschnitt~\ref{sec:implementation} deren Implementierung.

\subsection{SAT und die Sudoku-Kodierung} \label{sec:sat-encoding}

Die Solver in diesem Kapitel arbeiten auf aussagenlogischen Formeln in \emph{konjunktiver Normalform} (\acs{KNF})\acused{KNF}.
Damit ein \acs{SAT}-Solver ein Sudoku lösen kann, müssen dessen Regeln daher als \acs{KNF}-Formel vorliegen.
Abschnitt~\ref{sec:sat-cnf} führt die nötigen Grundbegriffe ein, Abschnitt~\ref{sec:encoding} übersetzt darauf aufbauend die Sudoku-Regeln in Klauseln.

\subsubsection{SAT und KNF} \label{sec:sat-cnf}

Ein \emph{Literal} ist eine boolesche Variable $x$ oder ihre Negation $\neg x$.
Eine \emph{Klausel} ist eine Disjunktion von Literalen. Sie ist erfüllt, sobald eines ihrer Literale wahr ist.
Eine \acs{KNF}-Formel ist eine Konjunktion von Klauseln \vgl[204]{davis1960}, zum Beispiel
\[
    F = (x_1 \lor \neg x_2) \land (x_2 \lor x_3) \land (\neg x_1 \lor \neg x_3).
\]
Eine Belegung mit Wahrheitswerten, die jede Klausel erfüllt, ist ein \emph{Modell} von $F$.
Hier ist $x_1 = x_2 = \text{wahr}$ und $x_3 = \text{falsch}$ ein solches.
\acs{SAT} ist die Frage, ob eine Formel ein Modell hat.

\acs{SAT} war das erste Problem, dessen NP-Vollständigkeit bewiesen wurde \vgl[151\psq]{cook1971}. Daher ist kein Algorithmus bekannt, der jede Instanz in Polynomialzeit löst.
Die beiden möglichen Antworten unterscheiden sich darin, wie schwer sie zu überprüfen sind.
Ein Modell lässt sich in linearer Zeit verifizieren, indem jede Klausel ausgewertet wird.
Für die Antwort \enquote{unerfüllbar} ist im Allgemeinen kein kurzes Zertifikat bekannt.
Die Beschränkung der Eingabe auf \acs{KNF} ist keine Einschränkung der Allgemeinheit: Jede aussagenlogische Formel lässt sich in eine erfüllbarkeitsäquivalente \acs{KNF}-Formel linearer Größe umformen \vgl[115\psq]{tseitin1970}.
Sudoku benötigt diesen Schritt nicht, da sich seine Regeln direkt als Klauseln formulieren lassen (Abschnitt~\ref{sec:encoding}).

Die Implementierung verwendet die ganzzahlige Notation des DIMACS-Formats, des gängigen Eingabeformats von \acs{SAT}-Solvern \vgl[2\psq]{dimacs1993}.
Die Variablen werden ab 1 nummeriert, das Literal $x_i$ wird als $i$ und $\neg x_i$ als $-i$ geschrieben, und eine Klausel ist eine Liste solcher ganzen Zahlen.
Aus $F$ wird damit \texttt{[[1, -2], [2, 3], [-1, -3]]}.

\subsubsection{Kodierung} \label{sec:encoding}

Die Kodierung folgt Lynce und Ouaknine~\parencite*[3\psqq]{lynce2006}.
Für jede Zelle $(r, c)$ und jede Ziffer $d$ gibt es eine boolesche Variable $x_{r,c,d}$, die genau dann wahr ist, wenn Zelle $(r, c)$ die Ziffer $d$ enthält.
Ein $n \times n$-Gitter benötigt daher $n^3$ Variablen, ein $9 \times 9$-Gitter also 729.
Die Implementierung nummeriert sie als
\[
    \mathrm{var}(r, c, d) = r \cdot n^2 + c \cdot n + d,
    \quad \text{mit } r, c \in \{0, \dots, n-1\} \text{ und } d \in \{1, \dots, n\}.
\]
Damit werden die Variablen lückenlos und eineindeutig auf $1, \dots, n^3$ abgebildet, wie es die DIMACS-Notation verlangt.
Das Dekodieren eines Modells ist die Umkehrrechnung zur Basis $n$.

Die Sudoku-Regeln werden zu vier Gruppen von Klauseln (Tabelle~\ref{tab:clause-groups}).
Zwei Gruppen fordern, dass etwas mindestens einmal gilt (\emph{\acl{ALO}}, \acs{ALO})\acused{ALO}, jeweils als eine einzige lange Klausel.
Die beiden anderen fordern, dass etwas höchstens einmal gilt (\emph{\acl{AMO}}, \acs{AMO})\acused{AMO}.
Sie werden paarweise kodiert: eine binäre Klausel für jedes Paar von Literalen, die nicht beide wahr sein dürfen.
Zusätzlich fügt jede vorgegebene Ziffer $g$ in Zelle $(r, c)$ die Unit-Klausel $(x_{r,c,g})$ hinzu.

\begin{table}[htb]
    \centering
    \caption{Klauselgruppen der Kodierung. $U$ ist eine Einheit: eine der $n$ Zeilen, $n$ Spalten oder $n$ Blöcke.}
    \label{tab:clause-groups}
    \small
    \renewcommand{\arraystretch}{1.3}
    \begin{tabularx}{\textwidth}{@{}>{\raggedright\arraybackslash}X l l r@{}}
        \toprule
        Gruppe & Klauseln & Anzahl & $9 \times 9$ \\
        \midrule
        \textbf{Zelle definiert} (\acs{ALO}): jede Zelle enthält mindestens eine Ziffer
            & $\bigwedge_{r,c} (x_{r,c,1} \lor \dots \lor x_{r,c,n})$
            & $n^2$ & 81 \\
        \textbf{Zelle eindeutig} (\acs{AMO}): jede Zelle enthält höchstens eine Ziffer
            & $\bigwedge_{r,c} \bigwedge_{d < d^{\prime}} (\neg x_{r,c,d} \lor \neg x_{r,c,d^{\prime}})$
            & $n^2 \binom{n}{2}$ & 2.916 \\
        \textbf{Einheit eindeutig} (\acs{AMO}): jede Ziffer kommt höchstens einmal pro Einheit vor
            & $\bigwedge_{U,d} \bigwedge_{(r,c) < (r^{\prime},c^{\prime}) \in U} (\neg x_{r,c,d} \lor \neg x_{r^{\prime},c^{\prime},d})$
            & $3n \cdot n \binom{n}{2}$ & 8.748 \\
        \textbf{Einheit definiert} (\acs{ALO}): jede Ziffer kommt mindestens einmal pro Einheit vor
            & $\bigwedge_{U,d} \bigvee_{(r,c) \in U} x_{r,c,d}$
            & $3n \cdot n$ & 243 \\
        \bottomrule
    \end{tabularx}
    \par\smallskip
    {\footnotesize Quelle: eigene Darstellung in Anlehnung an Lynce und Ouaknine~\parencite*[4\psq]{lynce2006}.}
\end{table}

Die \emph{minimale Kodierung} besteht nur aus den Gruppen \enquote{Zelle definiert} und \enquote{Einheit eindeutig} und ist bereits korrekt.
Eine Einheit hat $n$ Zellen, von denen jede mindestens eine Ziffer enthält, während jede der $n$ Ziffern höchstens einmal vorkommen darf.
Also enthält jede Zelle genau eine Ziffer (Schubfachprinzip).
Die beiden anderen Gruppen sind daher redundant.
Die \emph{erweiterte Kodierung} fügt sie dennoch hinzu: Sie ändern die Lösungsmenge nicht, geben dem Solver aber mehr zu propagieren.
So kann die Unit Propagation (Abschnitt~\ref{sec:dpll}) mithilfe von \enquote{Einheit definiert} eine Ziffer setzen, die nur in eine einzige Zelle einer Einheit passt. Menschliche Löser nennen dies ein \emph{Hidden Single}.

Für ein $9 \times 9$-Gitter hat die minimale Kodierung 8.829 Klauseln und die erweiterte 11.988, was mit Lynce und Ouaknine~\parencite*[4\psq]{lynce2006} übereinstimmt.
Beide wachsen mit $O(n^4)$.
Die zusätzlichen Klauseln zahlen sich aus, sobald der Solver suchen muss: Bei drei $16 \times 16$-Rätseln mit eindeutiger Lösung benötigte der \acs{CDCL}-Solver aus Abschnitt~\ref{sec:cdcl} im Median 3.134 Konflikte mit der minimalen und 71 mit der erweiterten Kodierung.

\subsection{Algorithmen} \label{sec:algorithms}

Beide Solver dieses Kapitels gehen auf den DPLL-Algorithmus zurück.
Abschnitt~\ref{sec:dpll} beschreibt DPLL, Abschnitt~\ref{sec:cdcl} dessen Erweiterung \acs{CDCL}, die aus Konflikten lernt.
Beide Algorithmen propagieren auf dieselbe Weise und unterscheiden sich nur im Umgang mit Konflikten; Abschnitt~\ref{sec:comparison} stellt sie gegenüber.

\subsubsection{DPLL} \label{sec:dpll}

Das erste Verfahren für \acs{SAT} eliminierte Variablen durch Resolution und benötigte in der Praxis zu viel Speicher \parencites[vgl.][210]{davis1960}[394]{davis1962}.
Die Ersetzung dieses Schritts durch Verzweigen mit Backtracking führte zum DPLL-Algorithmus, benannt nach Davis, Putnam, Logemann und Loveland \vgl[394]{davis1962}.

DPLL vereinfacht die Formel zunächst so weit wie möglich.
Die wichtigste Regel ist die \emph{Unit Propagation}: Sind alle Literale einer Klausel bis auf eines falsch, muss das verbleibende Literal wahr sein.
Dessen Belegung kann weitere Klauseln zu Unit-Klauseln machen, daher wird die Regel wiederholt, bis keine solche Klausel mehr übrig ist.
Optional belegt die \emph{Pure-Literal-Regel} zusätzlich jede Variable, die nur mit einem Vorzeichen vorkommt, so, dass alle ihre Vorkommen wahr sind.
Enthält danach eine Klausel nur noch falsche Literale, führt die aktuelle Belegung zu einem Konflikt.
Ist jede Klausel erfüllt, ist die Formel erfüllbar.
Andernfalls \emph{verzweigt} DPLL: Es wählt eine freie Variable $x$, setzt sie auf wahr und beginnt erneut mit der Propagation.
Führt dies zu einem Konflikt, kehrt der Solver per Backtracking zurück und versucht stattdessen $x = \text{falsch}$.
Scheitern beide Werte, wird der Konflikt an die vorherige Verzweigung weitergereicht \parencites[vgl.][209--212]{davis1960}[394\psq]{davis1962}.

Beide Vereinfachungsregeln erhalten die Erfüllbarkeit, und die Verzweigung probiert beide Werte jeder gewählten Variable aus.
Daher terminiert DPLL stets mit der korrekten Antwort; diese lautet \enquote{unerfüllbar}, sobald beide Werte der ersten gewählten Variable gescheitert sind.
Im schlechtesten Fall probiert es jedoch weiterhin alle $2^V$ Belegungen von $V$ Variablen aus.

Auf Sudoku angewandt führt die Unit Propagation die grundlegenden Züge menschlicher Löser aus.
Eine gesetzte Ziffer wird als Kandidat aus den übrigen Zellen ihrer Einheiten entfernt (\enquote{Einheit eindeutig}).
Eine Zelle mit nur noch einem Kandidaten, ein \emph{Naked Single}, wird über \enquote{Zelle definiert} gefüllt.
Ein \emph{Hidden Single}, also eine Ziffer, die nur in eine einzige Zelle einer Einheit passt, benötigt \enquote{Einheit definiert} und wird daher nur mit der erweiterten Kodierung gefunden.
Das Verzweigen entspricht dem Raten.

\subsubsection{CDCL} \label{sec:cdcl}

Gerät DPLL in einen Konflikt, macht es nur die letzte Entscheidung rückgängig und vergisst, warum der Konflikt entstanden ist.
Dieselbe Kombination von Belegungen kann daher in anderen Teilen der Suche immer wieder denselben Konflikt verursachen.
\acs{CDCL} analysiert stattdessen jeden Konflikt.
Es fügt die Ursache als neue Klausel zur Formel hinzu und springt direkt dorthin zurück, wo diese Ursache begann \vgl[509\psqq]{marquessilva1999}.
Zusammen mit den Techniken des Chaff-Solvers \vgl[530\psqq]{moskewicz2001} bildet dies die Grundlage praktisch aller modernen \acs{SAT}-Solver.

\acs{CDCL} behält die Unit Propagation von DPLL bei, speichert aber mehr Informationen.
Jede Entscheidung eröffnet eine neue \emph{Entscheidungsebene}.
Jedes danach propagierte Literal gehört zu dieser Ebene und merkt sich seinen \emph{Grund}: die Klausel, die zur Unit-Klausel wurde und es erzwungen hat.
Entscheidungen und Gründe bilden den \emph{Implikationsgraphen}: Seine Knoten sind die belegten Literale, und jedes propagierte Literal hat eingehende Kanten von den übrigen, falschen Literalen seines Grundes.
Ein Konflikt ist eine Klausel, deren Literale alle falsch sind.

Die Konfliktanalyse durchläuft diesen Graphen rückwärts, ausgehend von der Konfliktklausel.
In jedem Schritt resolviert sie die aktuelle Klausel mit dem Grund ihres zuletzt belegten Literals der aktuellen Ebene, wodurch dieses Literal durch seine Ursachen ersetzt wird.
Sie endet, sobald genau ein Literal der aktuellen Ebene übrig ist.
Dieses Literal ist der \emph{first \acl{UIP}} (first \acs{UIP})\acused{UIP}: die dem Konflikt nächstgelegene Belegung, durch die jeder Pfad von der aktuellen Entscheidung zum Konflikt verläuft \vgl[280\psqq]{zhang2001}.
Die resultierende Klausel enthält den negierten \acs{UIP} und ansonsten nur Literale früherer Ebenen.
Da sie durch Resolution hergeleitet wird, folgt sie aus der Formel und ändert deren Erfüllbarkeit nicht.
Sie schließt die Kombination von Belegungen aus, die den Konflikt verursacht hat.

Anschließend führt der Solver einen \emph{Backjump} zur zweithöchsten Entscheidungsebene der gelernten Klausel aus, bzw. zu Ebene 0, wenn die Klausel nur ein Literal enthält.
Auf dieser Ebene sind alle ihre Literale außer dem negierten \acs{UIP} falsch, sodass die Klausel eine Unit-Klausel ist und sofort eine neue Belegung erzwingt.
Ebenen, deren Entscheidungen für den Konflikt keine Rolle spielten, werden vollständig übersprungen, anders als beim chronologischen Backtracking von DPLL.
Da jede gelernte Klausel eine Belegung ausschließt, die bereits zu einem Konflikt geführt hat, terminiert die Suche weiterhin.

Vier weitere Techniken entscheiden darüber, wie gut dies in der Praxis funktioniert.
Die Heuristik \emph{\acl{VSIDS}} (\acs{VSIDS})\acused{VSIDS} wählt die Entscheidungsvariablen: Variablen, die an jüngsten Konflikten beteiligt waren, erhalten einen höheren Score, und alle Scores klingen mit der Zeit ab, sodass der Solver dort verzweigt, wo die letzten Konflikte auftraten \vgl[531\psq]{moskewicz2001}.
\emph{Restarts} setzen die Suche in Abständen, die durch die Luby-Folge vorgegeben sind, auf Ebene 0 zurück, behalten aber die gelernten Klauseln und die Scores bei, damit eine ungünstige frühe Entscheidung den Solver nicht festhält \parencites[vgl.][532]{moskewicz2001}[178]{luby1993}.
Beim \emph{Phase Saving} erhält eine Entscheidungsvariable den Wert, den sie zuletzt hatte, statt eines festen Werts; dies führt die Suche nach einem Restart zu vielversprechenden Belegungen zurück \vgl[297]{pipatsrisawat2007}.
Schließlich kosten gelernte Klauseln Speicher und verlangsamen die Propagation, weshalb der Solver regelmäßig die weniger nützliche Hälfte von ihnen löscht.
Die Nützlichkeit wird über die \emph{\acl{LBD}} (\acs{LBD})\acused{LBD} gemessen, die Anzahl verschiedener Entscheidungsebenen in einer Klausel. Klauseln mit niedriger \acs{LBD} werden behalten \vgl[401\psq]{audemard2009}.

\subsubsection{DPLL vs.\ CDCL} \label{sec:comparison}

Tabelle~\ref{tab:dpll-vs-cdcl} fasst die Unterschiede zusammen.

\begin{table}[htb]
    \centering
    \caption{DPLL und \acs{CDCL} im Vergleich}
    \label{tab:dpll-vs-cdcl}
    \small
    \renewcommand{\arraystretch}{1.3}
    \begin{tabularx}{\textwidth}{@{}>{\raggedright\arraybackslash}p{3.6cm} >{\raggedright\arraybackslash}X >{\raggedright\arraybackslash}X@{}}
        \toprule
         & DPLL & \acs{CDCL} \\
        \midrule
        Reaktion auf einen Konflikt
            & versucht den entgegengesetzten Wert der letzten offenen Entscheidung
            & lernt eine Klausel, die durch Propagation eine neue Belegung erzwingt \\
        Backtracking
            & chronologisch, jeweils eine Ebene
            & nicht chronologisch, zur zweithöchsten Ebene der gelernten Klausel \\
        Gedächtnis für frühere Konflikte
            & keines
            & gelernte Klauseln, nach \acs{LBD} ausgedünnt \\
        Reihenfolge der Entscheidungen
            & Verzweigungsheuristik, unabhängig von Konflikten
            & \leavevmode\acs{VSIDS}, folgt jüngsten Konflikten \\
        Restarts
            & nein
            & ja, gelernte Klauseln bleiben erhalten \\
        \bottomrule
    \end{tabularx}
    \par\smallskip
    {\footnotesize Quelle: eigene Darstellung.}
\end{table}

Der Unterschied zeigt sich erst, wenn der Solver suchen muss.
Bei 20~$9 \times 9$-Rätseln mit eindeutiger Lösung und der minimalen Kodierung mussten beide Solver kaum suchen: DPLL benötigte im Median 24 Konflikte und 0,03\,s, \acs{CDCL} 9 Konflikte und 0,02\,s.
Von drei $16 \times 16$-Rätseln mit eindeutiger Lösung löste DPLL dagegen keines innerhalb des Zeitlimits von 20\,s, während \acs{CDCL} alle drei in einem Median von 1,3\,s löste.

\subsection{Implementierung} \label{sec:implementation}

Die beiden \acs{SAT}-Solver, DPLL und \acs{CDCL}, wurden ohne \acs{SAT}-Bibliothek in Python implementiert.
Abschnitt~\ref{sec:architecture} beschreibt die gemeinsame Architektur, die Abschnitte~\ref{sec:dpll-impl} und~\ref{sec:cdcl-impl} die Datenstrukturen der beiden Solver und Abschnitt~\ref{sec:trace} das Suchprotokoll, auf dem die Visualisierung aufbaut.

\subsubsection{Architektur} \label{sec:architecture}

Die Implementierung folgt den drei Schritten aus der Einleitung dieses Kapitels.
\texttt{encode} wandelt das Gitter in eine Liste von Klauseln und die Anzahl der Variablen um (Abschnitt~\ref{sec:encoding}).
Anschließend sucht einer der \acs{SAT}-Solver ein Modell, und \texttt{decode\_model} wandelt die wahren Variablen dieses Modells zurück in ein ausgefülltes Gitter.
Findet der Solver kein Modell, hat das Rätsel keine Lösung.

Alle \acs{SAT}-Solver teilen dieselbe Schnittstelle, \texttt{solve\_cnf(clauses, num\_vars)}.
Sie gibt ein Modell zurück, das für jede Variable genau ein Literal, $v$ oder $-v$, enthält, oder \texttt{None}, falls die Formel unerfüllbar ist, und lässt die Klauseln unverändert.
Zwei optionale Argumente sind allen Solvern gemeinsam: \texttt{stats} erfasst die Anzahl der Entscheidungen, Propagationen und Konflikte, und \texttt{trace} zeichnet jeden Suchschritt auf (Abschnitt~\ref{sec:trace}).
Das Aufzeichnen eines Protokolls verlangsamt einen Solver um 30\,\% bis hin zum Faktor 2,3.
Ohne Protokoll wird der Aufzeichnungscode übersprungen, und die verbleibenden Prüfungen kosten weniger als 5\,\% der Laufzeit, was innerhalb des Messrauschens liegt.
Dank dieser Schnittstelle lassen sich die Solver beliebig austauschen und an jeder \acs{KNF}-Formel testen, nicht nur an Sudoku.

Nichts im Code hängt von einer festen Gittergröße ab.
Jede Funktion leitet $n$ aus dem Gitter und die Blockgröße $b = \sqrt{n}$ aus $n$ ab, sodass derselbe Code $4 \times 4$-, $9 \times 9$- und $16 \times 16$-Gitter verarbeitet.

\subsubsection{Vom einfachen zum effizienten DPLL} \label{sec:dpll-impl}

Die erste DPLL-Implementierung folgt Abschnitt~\ref{sec:dpll} wörtlich.
Sie ist rekursiv, und jede Belegung erzeugt eine neue, vereinfachte Kopie der Formel ohne die erfüllten Klauseln und die falschen Literale.
Dadurch ist sie leicht zu überprüfen, doch jeder Suchschritt kostet Zeit proportional zur Größe der gesamten Formel.

Die effiziente Version führt dieselbe Suche mit anderen Datenstrukturen durch.
Statt die Formel zu kopieren, führt sie einen \emph{Trail}: die Liste der belegten Literale in der Reihenfolge, in der sie gesetzt wurden.
Backtracking kürzt den Trail einfach auf die Länge, die er vor der gescheiterten Entscheidung hatte.
Die offenen Entscheidungen liegen auf einem separaten Stack, sodass die Suche iterativ ist und nicht durch die Rekursionstiefe von Python begrenzt wird.

Die Unit Propagation verwendet \emph{Two Watched Literals} \vgl[531]{moskewicz2001}.
Jede Klausel beobachtet zwei ihrer Literale, die nicht falsch sind, und wird nur betrachtet, wenn eines davon falsch wird.
Der Solver sucht dann in der Klausel nach einem anderen nicht falschen Literal, das stattdessen beobachtet wird.
Gibt es keines, ist die Klausel zur Unit-Klausel oder zum Konflikt geworden.
Da Backtracking Belegungen nur aufhebt, bleiben die beobachteten Literale gültig und müssen nicht aktualisiert werden.

Entscheidungen folgen der Heuristik \texttt{shortest}: dem ersten freien Literal der kürzesten offenen Klausel, die nur positive Literale enthält.
Auf Sudoku ist dies die Zelle oder Einheit mit den wenigsten verbleibenden Kandidaten.
Die Beschränkung auf positive Klauseln ist wichtig: Die offenen binären \acs{AMO}-Klauseln sind stets die kürzesten, und eine Entscheidung \enquote{$x$ ist falsch} bewirkt wenig.

Sobald der Solver suchen muss, sparen diese Änderungen ein bis zwei Größenordnungen ein.
Auf dem leeren $16 \times 16$-Gitter mit der minimalen Kodierung benötigte die einfache Version 15,9\,s und die effiziente 0,36\,s.
\acs{CDCL} (Abschnitt~\ref{sec:cdcl-impl}) baut auf demselben Trail und denselben Watched Literals auf.

\subsubsection{CDCL} \label{sec:cdcl-impl}

Der \acs{CDCL}-Solver erweitert den effizienten DPLL.
Für jede Variable speichert er zusätzlich die Entscheidungsebene und den Grund, also die Klausel, die sie erzwungen hat. Entscheidungen haben keinen Grund.
Die Konfliktanalyse ist ein einziger Durchlauf rückwärts entlang des Trails (Listing~\ref{lst:analyze}).
Sie markiert die Variablen der Konfliktklausel und zählt, wie viele davon zur aktuellen Ebene gehören.
Anschließend nimmt sie das zuletzt belegte markierte Literal vom Trail und fährt, solange noch mehr als ein Literal der aktuellen Ebene offen ist, mit dem Grund dieses Literals fort.
Ist nur noch eines übrig, ist das gerade entnommene Literal der first \acs{UIP}.
Literale früherer Ebenen gehen direkt in die gelernte Klausel ein.
Die Funktion setzt den negierten \acs{UIP} an die erste und ein Literal der Backjump-Ebene an die zweite Stelle, damit diese beiden sofort beobachtet werden können, und berechnet im selben Durchlauf die \acs{LBD}.

% Pseudocode als Listing im Fließtext (nicht gleitend, Eintrag im Listingverzeichnis)
\par\medskip
\noindent\begin{minipage}{\linewidth}
\fcolorbox{black}{boxgray}{%
\begin{minipage}{\dimexpr\linewidth-2\fboxsep-2\fboxrule\relax}
\small
\begin{algorithmic}[1]
\Require Konfliktklausel $K_\bot$, deren Literale alle falsch sind
\Ensure gelernte Klausel $L$, Backjump-Ebene $b$, \acs{LBD} $\ell$
\Function{Analysiere}{$K_\bot$}
  \State $e \gets$ aktuelle Entscheidungsebene
  \State $L \gets [\,\cdot\,]$ \Comment{Platz 0 für das erzwingende Literal}
  \State $o \gets 0$ \Comment{offene Literale der Ebene $e$}
  \State $i \gets |T| - 1$ \Comment{$T$: Trail}
  \State $K \gets K_\bot$; \; $p \gets \text{nil}$
  \Repeat
    \ForAll{$q \in K$ mit $q \neq p$}
      \State $v \gets \operatorname{var}(q)$
      \If{$v$ nicht markiert \textbf{und} $\operatorname{ebene}(v) > 0$}
        \State markiere $v$; \; \Call{ErhöheScore}{$v$} \Comment{VSIDS}
        \If{$\operatorname{ebene}(v) = e$}
          \State $o \gets o + 1$
        \Else
          \State hänge $q$ an $L$ an
        \EndIf
      \EndIf
    \EndFor
    \While{$\operatorname{var}(T[i])$ nicht markiert} \Comment{zuletzt markiertes Literal}
      \State $i \gets i - 1$
    \EndWhile
    \State $p \gets T[i]$; \; $i \gets i - 1$
    \State entferne Markierung von $\operatorname{var}(p)$; \; $o \gets o - 1$
    \If{$o > 0$}
      \State $K \gets \operatorname{grund}(\operatorname{var}(p))$ \Comment{mit dem Grund von $p$ resolvieren}
    \EndIf
  \Until{$o = 0$} \Comment{$p$ ist der first \acs{UIP}}
  \State $L[0] \gets \neg p$
  \State entferne Markierungen aller Variablen in $L[1..]$
  \State $b \gets 0$
  \If{$|L| > 1$}
    \State $j \gets$ Index in $L[1..]$ mit höchster Entscheidungsebene \Comment{Backjump-Ebene}
    \State vertausche $L[1]$ und $L[j]$
    \State $b \gets \operatorname{ebene}(\operatorname{var}(L[1]))$
  \EndIf
  \State $\ell \gets |\{\operatorname{ebene}(\operatorname{var}(q)) : q \in L[1..]\}| + 1$
  \State \Return $(L, b, \ell)$
\EndFunction
\end{algorithmic}
\end{minipage}}
\captionof{lstlisting}{First-UIP-Konfliktanalyse als Pseudocode (nach \texttt{analyze} in \texttt{sat/cdcl.py})}
\label{lst:analyze}
{\footnotesize Quelle: eigene Darstellung.}
\end{minipage}
\par\medskip

Nach dem Backjump wird die gelernte Klausel zum Grund des neu erzwungenen Literals.
\acs{VSIDS} erhöht den Score jeder während der Analyse markierten Variable.
Statt nach jedem Konflikt alle Scores zu verringern, erhöht der Solver das Inkrement um den Faktor $1/0{,}95$, was dieselbe Wirkung hat.
Freie Variablen werden einem Heap entnommen, aus dem veraltete Einträge erst entfernt werden, wenn sie an die Reihe kommen.
Restarts erfolgen nach $100 \cdot \mathrm{luby}(i)$ Konflikten.
Ab 2.000 Konflikten löscht der Solver in wachsenden Abständen die Hälfte der gelernten Klauseln mit der höchsten \acs{LBD}, behält aber Klauseln mit einer \acs{LBD} von höchstens 2 sowie Klauseln, die gerade als Grund dienen.
Jede dieser Techniken lässt sich einzeln abschalten, sodass ihre Wirkung gemessen werden kann.

\subsubsection{Suchprotokoll} \label{sec:trace}

Um die Suche sichtbar zu machen, können beide Solver ein \emph{Protokoll} (\emph{Trace}) aufzeichnen: eine Liste mit einem Eintrag für jeden Suchschritt, die die Visualisierung anschließend Schritt für Schritt abspielen kann.
Eine Propagation wird mit dem Literal, seiner Entscheidungsebene und seinem Grund protokolliert.
Weitere Einträge kennzeichnen Entscheidungen, Konflikte und Backtracking, bei \acs{CDCL} zudem gelernte Klauseln, Restarts und das Löschen von Klauseln.

Auf Sudoku übertragen entspricht jeder Eintrag direkt einer Änderung des Gitters.
Eine auf wahr gesetzte Variable $x_{r,c,d}$ trägt die Ziffer $d$ in Zelle $(r, c)$ ein, eine auf falsch gesetzte entfernt $d$ aus den Kandidaten dieser Zelle.
Die meisten Schritte sind von der zweiten Art: In einem \acs{CDCL}-Protokoll des schweren $9 \times 9$-Rätsels \enquote{AI Escargot} entfernen 82\,\% aller Einträge einen Kandidaten.
Ein \acs{SAT}-Solver lässt sich daher am besten auf einem Gitter mit Kandidaten darstellen, nicht als Folge gesetzter Ziffern.
Mit der minimalen Kodierung umfasst das Protokoll von AI Escargot 1.383 Einträge für DPLL und 11.449 für \acs{CDCL}.

\providecommand{\plotdir}{../evaluation/plots/}

\section{Evaluierung}\label{sec:evaluation}

In diesem Kapitel werden Backtracking, DPLL und CDCL anhand eigener Messungen auf 9x9-Sudokus verglichen. Betrachtet werden die Laufzeit, der Suchaufwand und die Verlässlichkeit der Verfahren bei steigender Schwierigkeit sowie bei unlösbaren und besonders schweren Rätseln. In Abschnitt~\ref{sec:versuchsaufbau} werden Rätsel, Messung und statistische Auswertung beschrieben. Abschnitt~\ref{sec:ergebnisse} enthält die Ergebnisse, in Abschnitt~\ref{sec:diskussion} werden sie eingeordnet.

\subsection{Versuchsaufbau}\label{sec:versuchsaufbau}
Mit der Evaluierung sollen vier Fragen beantwortet werden. Wie ändern sich die Laufzeiten von Backtracking, DPLL und CDCL, wenn mehr Felder leer sind? Was kostet die Umwandlung in eine SAT-Formel? Welche Bausteine der SAT-Solver bringen etwas? Wie verhalten sich die Verfahren bei unlösbaren und besonders schweren Rätseln?

Alle Messungen liefen nacheinander auf einem Apple M2 mit acht Kernen unter macOS mit CPython 3.14.3 und py-sudoku 2.0.0. Alle drei Verfahren sind in reinem Python geschrieben. Die Messung importiert den Projektcode und verändert ihn nicht.

Die Schwierigkeit eines Rätsels ist der Anteil leerer Felder. Es gibt sieben Stufen von 30 bis 70 \%. Die Familie \emph{eindeutig} enthält Rätsel mit genau einer Lösung. Sie entstehen, indem aus einer vollständigen Lösung zufällig Felder entfernt werden, solange die Lösung eindeutig bleibt. Dadurch hat die Stufe 70 \% im Mittel nur 55,9 statt 56 leere Felder. Die Familie \emph{zufall} wird wie in der Oberfläche des Projekts erzeugt, ihre Rätsel haben oft mehrere Lösungen, bei 30 \% leeren Feldern in 10 \% und bei 70 \% in allen Rätseln. Der Hauptvergleich nutzt \emph{eindeutig}, \emph{zufall} dient als Gegenprobe. Pro Stufe und Familie gibt es 100 Rätsel, und alle Verfahren lösen dieselben Rätsel. Weitere Experimente verwenden kleinere Mengen von je 50 Rätseln bei 50, 60 und 70 \% leeren Feldern (Heuristiken, Konfigurationen, Kodierungen), unlösbare Rätsel und sechs handverlesene schwere Rätsel.

Jedes Verfahren darf höchstens 20 Sekunden Wanduhrzeit je Rätsel brauchen, sonst wird der Lauf abgebrochen. Ein Abbruch zählt in Median und Perzentilen mit 20 s und im PAR-2-Wert mit 40 s. Die Zeit eines Rätsels ist der Median aus bis zu fünf Läufen, Läufe über 2 s werden nicht wiederholt. Während der Zeitnahme ist die Garbage Collection ausgeschaltet, bei SAT zählt das Erzeugen der Formel mit. Die zusätzlich erfasste CPU-Zeit weicht bei der Hauptmessung bei 90 \% der Läufe um weniger als 0,2 \% von der Wanduhrzeit ab. Diagramme und Tests beruhen auf der Wanduhrzeit. Speicher wurde nicht gemessen.

Je Stufe werden Median mit 95-\%-Bootstrap-Intervall (2000 Wiederholungen), 90. Perzentil, Maximum, Abbruchquote und PAR-2 berichtet. Zwei Verfahren werden gepaart verglichen, und zwar über das geometrische Mittel der Quotienten der Zeiten mit Bootstrap-Intervall und einen Wilcoxon-Test mit Normalapproximation. Das Niveau beträgt 5 \%, korrigiert nach Bonferroni über alle Tests einer Tabelle. Aufwandszähler gehören zu verschiedenen Verfahren und werden nicht auf eine Achse gelegt. Jede Lösung wird auf Gültigkeit geprüft.

\subsection{Ergebnisse}\label{sec:ergebnisse}
\paragraph{Hauptvergleich.} Tabelle~\ref{tab:e1} zeigt die Laufzeiten auf \emph{eindeutig}, Abbildung~\ref{abb:e1} den Verlauf. Das Backtracking startet bei 30 \% leeren Feldern mit einem Median von 0,19 ms und steigt bei 70 \% auf 126 ms, also auf das 660-Fache. Die SAT-Verfahren bleiben fast flach. CDCL braucht 7,17 ms bei 30 \% und 8,37 ms bei 70 \%, DPLL 8,82 ms und 9,50 ms. Bis 55 \% ist das Backtracking schneller, bei 30 \% etwa 37-mal und bei 55 \% noch knapp 5-mal gegenüber CDCL. Bei 60 \% liegt sein Median mit 4,65 ms unter dem von CDCL mit 7,18 ms, sein 90. Perzentil mit 25,4 ms aber deutlich darüber. Auch der Mittelwert, der hier dem PAR-2-Wert entspricht, ist mit 11,6 ms größer als bei CDCL mit 7,34 ms. Bei 60 \% ist der Unterschied nicht signifikant. Ab 65 \% ist CDCL schneller, um den Faktor 3,2 bei 65 \% und 13,5 bei 70 \%. DPLL liegt in allen Stufen 15 bis 23 \% hinter CDCL.

\begin{table}[htb]
\centering\small\setlength{\tabcolsep}{3.5pt}
\caption{Laufzeit in ms auf \emph{eindeutig}. Quotient: geometrisches Mittel der Zeit von Backtracking geteilt durch die von CDCL. Test: Wilcoxon, Backtracking gegen CDCL, Bonferroni-korrigiert. Ohne Entscheidung: Anteil der Rätsel (\%), die CDCL ohne Entscheidung löst.}\label{tab:e1}
\begin{tabular}{@{}r r r r r r l c r@{}}
\toprule
\shortstack[r]{Leer\\(\%)} & \shortstack[r]{BT\\Median} & \shortstack[r]{BT\\p90} & \shortstack[r]{DPLL\\Median} & \shortstack[r]{CDCL\\Median} & Quotient & \shortstack[l]{95-\%-\\Intervall} & \shortstack{signi-\\fikant} & \shortstack[r]{ohne\\Entsch.} \\
\midrule
30 & 0,19 & 0,23 & 8,82 & 7,17 & 0,027 & 0,026 bis 0,028 & ja & 100 \\
40 & 0,29 & 0,51 & 8,79 & 7,16 & 0,045 & 0,043 bis 0,049 & ja & 100 \\
50 & 0,63 & 2,63 & 8,79 & 7,17 & 0,12 & 0,10 bis 0,14 & ja & 100 \\
55 & 1,29 & 4,75 & 8,81 & 7,18 & 0,21 & 0,17 bis 0,25 & ja & 89 \\
60 & 4,65 & 25,4 & 8,81 & 7,18 & 0,64 & 0,50 bis 0,83 & nein & 65 \\
65 & 31,5 & 175 & 8,90 & 7,50 & 3,24 & 2,33 bis 4,48 & ja & 15 \\
70 & 126 & 1\,146 & 9,50 & 8,37 & 13,5 & 9,57 bis 19,6 & ja & 0 \\
\bottomrule
\end{tabular}
\end{table}

\begin{figure}[htb]
\centering
\includegraphics[width=0.8\textwidth]{\plotdir e1_laufzeit_eindeutig}
\caption{Median der Laufzeit je Stufe (\emph{eindeutig}), Band: 25. bis 75. Perzentil}\label{abb:e1}
\end{figure}

Die Zeiten des Backtrackings streuen stark. Bei 70 \% liegt der Median bei 126 ms, das 90. Perzentil bei 1,1 s und die langsamste Zeit bei 14,9 s. Bei CDCL ist die langsamste Zeit mit 19,8 ms nur das 2,4-Fache des Medians. Auf \emph{eindeutig} gibt es keinen Abbruch. Auf \emph{zufall} gibt es bei der Hauptmessung genau einen, beim Backtracking bei 70 \%. Dort wird das Backtracking später langsam, bei 70 \% beträgt sein Median 18,8 ms, und CDCL ist nur 2,7-mal statt 13,5-mal schneller.

\paragraph{Suchaufwand und Kosten der Formel.} Der Median der Schritte des Backtrackings steigt von 26 bei 24 leeren Feldern auf 99 bei 40 und auf 20\,830 bei etwa 56. Ob ein SAT-Solver suchen muss, hängt ebenfalls von der Stufe ab (Abbildung~\ref{abb:ohne}). Bis 50 \% löst die Propagation allein alle Rätsel, bei 55 \% sind es 89 \%, bei 60 \% 65 \%, bei 65 \% 15 \% und bei 70 \% kein Rätsel. Bei \emph{zufall} ist der Anteil kleiner, bei 50 \% nur 39 \%. Er entspricht fast dem Anteil der Rätsel mit genau einer Lösung (40 \%). Bei mehreren Lösungen bleibt mindestens ein Feld unbestimmt, das die Propagation nicht festlegen kann, also muss geraten werden. Die Schwierigkeit hängt nur teilweise vom Verfahren ab. Über alle Stufen korrelieren die Schritte des Backtrackings mit den Konflikten von CDCL mit $\rho = 0{,}70$, innerhalb einer Stufe ab 55 \% nur mit 0,05 bis 0,47. Beide Zähler wachsen also vor allem gemeinsam mit der Zahl leerer Felder. Die Formel hat bei jedem Rätsel 729 Variablen und 8829 Klauseln, dazu kommt eine Klausel je Vorgabe. Das Erzeugen dauert im Median 1,08 ms, das sind bei CDCL 13 bis 15 \% und bei DPLL 11,5 bis 12,3 \% der Gesamtzeit.

\begin{figure}[htb]
\centering
\includegraphics[width=0.8\textwidth]{\plotdir e2_ohne_suche_eindeutig}
\caption{Anteil der Rätsel, die ohne Rückschritt (Backtracking) bzw. ohne Entscheidung (SAT) gelöst werden (\emph{eindeutig})}\label{abb:ohne}
\end{figure}

\paragraph{Bausteine der SAT-Solver.} Tabelle~\ref{tab:bausteine} zeigt den Einfluss einzelner Bausteine auf Zeit und Konflikte. Bei 50 \% braucht kein DPLL-Lauf eine Entscheidung, bei 60 \% ist der Median der Entscheidungen null. Die Medianzeiten unterscheiden sich trotzdem (DPLL: \emph{first} 6,9 ms, \emph{shortest\_any} 7,4 ms, \emph{shortest} 8,8 ms, \emph{jw} 8,9 ms). Der Unterschied kommt also nicht von der Suche, sondern von festen Kosten des jeweiligen Auswahlverfahrens. Bei 70 \% zählt die Qualität der Auswahl. Der Median der Konflikte liegt bei \emph{shortest} und \emph{shortest\_any} bei 13, bei \emph{jw} bei 18,5 und bei \emph{first} bei 42,5. Das Maximum sinkt von 958 (\emph{first}) auf 162 (\emph{shortest}).

\begin{table}[htb]
\centering\small
\caption{Verhältnis zur Standardkonfiguration als geometrisches Mittel über 150 Paare (3 Stufen mit je 50 Rätseln) mit 95-\%-Intervall. Bei den Konflikten wurde 1 addiert, weil viele Läufe keinen Konflikt haben. Bei der erweiterten Kodierung ist die Referenz die minimale Kodierung. Alle drei Bausteine aus: ohne Phase-Saving, Neustarts und Klausellöschung.}\label{tab:bausteine}
\begin{tabular}{@{}l l l@{}}
\toprule
Konfiguration & Zeit & Konflikte \\
\midrule
CDCL Heuristik shortest statt VSIDS & 1,20 (1,18 bis 1,21) & 0,85 (0,78 bis 0,94) \\
CDCL ohne Phase-Saving & 0,99 (0,97 bis 1,00) & 0,93 (0,87 bis 0,99) \\
CDCL ohne Neustarts & 1,00 (1,00 bis 1,00) & 1,00 (1,00 bis 1,00) \\
CDCL ohne Klausellöschung & 1,00 (1,00 bis 1,00) & 1,00 (1,00 bis 1,00) \\
CDCL Anfangsphase wahr & 1,00 (0,98 bis 1,02) & 1,02 (0,94 bis 1,12) \\
CDCL alle drei Bausteine aus & 0,99 (0,97 bis 1,00) & 0,93 (0,87 bis 0,99) \\
DPLL mit erweiterter Kodierung & 1,33 (1,30 bis 1,35) & 0,50 (0,41 bis 0,59) \\
CDCL mit erweiterter Kodierung & 1,28 (1,25 bis 1,31) & 0,49 (0,42 bis 0,57) \\
\bottomrule
\end{tabular}
\end{table}

Bei CDCL verändert nur die Entscheidungsheuristik die Zeit. Mit \emph{shortest} statt VSIDS sinken Konflikte (Faktor 0,85) und Entscheidungen (0,78), die Zeit steigt aber um den Faktor 1,20. Phase-Saving abzuschalten spart etwa 1 \% Zeit. Neustarts und Klausellöschung haben in dieser Messreihe keinen Effekt, denn sie werden nie ausgelöst. Neustarts beginnen nach 100 Konflikten, die erste Löschung nach 2000, der größte Wert in dieser Messreihe sind 87 Konflikte. Selbst bei den schweren Rätseln sind es höchstens 534 Konflikte und vier Neustarts, eine Löschung kommt nie vor. CDCL mit \emph{shortest} (8,85, 8,90 und 9,63 ms) und DPLL mit \emph{shortest} (8,80, 8,81 und 9,55 ms) sind bei 50, 60 und 70 \% praktisch gleich schnell. Die erweiterte Kodierung hat 36 \% mehr Klauseln und halbiert die Konflikte (Faktor 0,49 bei CDCL, 0,50 bei DPLL). Die Zeit steigt dennoch um 28 \% (CDCL) bzw. 33 \% (DPLL).

\begin{figure}[htb]
\centering
\includegraphics[width=0.8\textwidth]{\plotdir e5_cdcl_ablation}
\caption{Laufzeit der CDCL-Konfigurationen}\label{abb:e5}
\end{figure}

\paragraph{Unlösbare und schwere Rätsel.} Bei den unlösbaren Rätseln stellen beide SAT-Verfahren die Unlösbarkeit bei allen 51 Rätseln allein durch Propagation fest, ohne Entscheidung und mit genau einem Konflikt. Der Median liegt bei 6,73 ms (CDCL) und 8,08 ms (DPLL). Das Backtracking hat einen Median von 0,55 ms, ist aber bei dem handverlesenen unlösbaren Rätsel mit 72 leeren Feldern nach 20 s abgebrochen worden. Sein PAR-2-Wert liegt deshalb bei 785 ms. Tabelle~\ref{tab:e8} zeigt die schweren Rätsel. Backtracking wird bei zwei der sechs Rätsel abgebrochen. CDCL löst alle in höchstens 80 ms, DPLL in höchstens 318 ms. Die PAR-2-Werte betragen 13\,450 ms (Backtracking), 86,9 ms (DPLL) und 28,7 ms (CDCL). Kein Verfahren gewinnt überall. CDCL ist bei drei Rätseln am schnellsten, DPLL bei zwei, Backtracking bei einem. Wegen der sechs Rätsel gibt es hier keine Tests.

\begin{table}[htb]
\centering\small
\caption{Laufzeit in ms bei handverlesenen schweren Rätseln, schnellster Wert fett. Abbruch nach 20\,000 ms.}\label{tab:e8}
\begin{tabular}{@{}l r r r r@{}}
\toprule
Rätsel & leer & Backtracking & DPLL & CDCL \\
\midrule
Inkala 2012 & 60 & 502,1 & 95,1 & \textbf{79,7} \\
AI Escargot & 58 & 91,0 & \textbf{9,4} & 22,2 \\
Anti-Backtracking & 64 & $\geq$ 20\,000 & 318,4 & \textbf{15,2} \\
Norvig hard1 & 64 & $\geq$ 20\,000 & \textbf{12,0} & 15,1 \\
ungünstig für SAT (viele Konflikte) & 57 & 68,2 & 58,9 & \textbf{27,3} \\
leicht für Backtracking, schwer für SAT & 57 & \textbf{9,9} & 27,8 & 12,9 \\
\bottomrule
\end{tabular}
\end{table}

Von 6621 Läufen wurden vier abgebrochen, jeweils beim Backtracking (einmal bei der Hauptmessung, einmal bei den unlösbaren und zweimal bei den schweren Rätseln). Alle beendeten Läufe waren korrekt, und es gab keinen Widerspruch zwischen den Verfahren.

\subsection{Diskussion}\label{sec:diskussion}
Kein Verfahren ist durchweg besser. Auf \emph{eindeutig} ist das Backtracking bis 55 \% leeren Feldern schneller, ab 65 \% CDCL, und dazwischen gibt es bei 60 \% keinen sicheren Unterschied. Der Grund liegt in der Form der Kosten. Die Formel ist für jedes 9x9-Rätsel gleich groß, und bei wenigen leeren Feldern löst die Propagation fast alles. Die Zeit der SAT-Verfahren besteht deshalb zum größten Teil aus festen Kosten von 7 bis 9 ms, von denen das Erzeugen der Formel nur etwa 1 ms ausmacht. Das Backtracking hat fast keine festen Kosten, dafür wächst sein Aufwand mit der Suche stark. Die Kurven müssen sich dort kreuzen, wo die Suche des Backtrackings so teuer wird wie diese festen Kosten. Das passt dazu, dass der Median des Backtrackings bei 60 \% mit 4,65 ms noch knapp unter und bei 65 \% mit 31,5 ms klar über diesen Kosten liegt. Die Lage des Kreuzungspunkts hängt von den Umsetzungen ab, die Form der Kurven nicht.

Wichtiger als der Median ist die Verlässlichkeit. Bei 60 \% hat das Backtracking den besseren Median, aber einen schlechteren Mittelwert, weil wenige sehr lange Läufe den Mittelwert bestimmen. Bei 70 \% ist sein langsamster Lauf rund 120-mal so lang wie sein Median, bei CDCL nur 2,4-mal. Bei unlösbaren und schweren Rätseln ist es ebenso. Die SAT-Verfahren brauchen für alle unlösbaren Rätsel etwa gleich lang, das Backtracking wurde beim schwersten Fall abgebrochen. Wer eine Obergrenze für die Wartezeit braucht, ist mit SAT besser beraten.

Der erwartete Vorteil von CDCL gegenüber DPLL ist in diesen Daten klein und hat eine andere Ursache als das Lernen. DPLL ist je Stufe 15 bis 23 \% langsamer. Dieselbe Lücke entsteht, wenn man bei CDCL die Heuristik durch die von DPLL ersetzt (Faktor 1,20). CDCL mit \emph{shortest} und DPLL sind praktisch gleich schnell. Das Lernen senkt zwar die Konflikte, bei 70 \% im Median von 13 auf 6,5, spart aber keine messbare Zeit. Neustarts und Löschung werden kaum ausgelöst, weil sie erst nach 100 bzw. 2000 Konflikten beginnen und die Rätsel meist weniger erzeugen. Man kann daraus nur schließen, dass diese Bausteine bei 9x9-Sudokus hier nicht gebraucht werden. Über größere Gitter sagen die Daten nichts.

Bei der erweiterten Kodierung bedeutet weniger Suche nicht automatisch weniger Zeit. Sie halbiert die Konflikte, kostet aber 36 \% mehr Klauseln und damit 28 bis 33 \% mehr Zeit. Konflikte sind bei Sudoku so selten und billig, dass zusätzliche Klauseln mehr kosten, als sie sparen. Dasselbe gilt für die Heuristik \emph{shortest} bei CDCL.

Rätsel mit mehreren Lösungen verhalten sich anders als eindeutige. Sie erzwingen im SAT-Solver mindestens eine Entscheidung, sind für das Backtracking aber leichter. Warum, wurde nicht untersucht. Wegen dieses Unterschieds ist die Familie \emph{eindeutig} der ehrlichere Vergleich. Weil außerdem die Schwierigkeit innerhalb einer Stufe zwischen den Verfahren nur schwach zusammenhängt, sagt der Anteil leerer Felder den Aufwand im Mittel voraus, aber nicht für das einzelne Rätsel. Die handverlesenen Rätsel zeigen das, denn zwei sind für Backtracking sehr schwer und für SAT leicht, bei einem ist es umgekehrt.

Rahmen und Einschränkungen: Alle Verfahren sind in Python geschrieben, absolute Zeiten und Quotienten gelten nur für diese Umsetzungen und diesen Rechner. Alle Rätsel stammen aus demselben Generator, die Aussagen des Wilcoxon-Vorzeichen-Rang-Tests mit Bonferroni-Korrektur und der Bootstrap-Intervalle beziehen sich auf diese Rätselmengen. Es wurden nur 9x9-Sudokus gelöst, Speicherbedarf und andere Gittergrößen wurden nicht untersucht.



	% Anhang derzeit ohne Inhalt und daher deaktiviert.
	% Zum Aktivieren die folgenden Zeilen einkommentieren und appendix/index.tex befüllen.

	% Anhang-Nummerierung
	% \renewcommand{\thetable}{\Alph{section}.\arabic{table}}
	% \renewcommand{\thefigure}{\Alph{section}.\arabic{figure}}
	% \renewcommand{\thelstlisting}{\Alph{section}.\arabic{lstlisting}}

	% Abschluss: Reihenfolge gem. Richtlinie V 2.0 (Anhang VOR Literatur)

    % \begin{appendix}
	% 	%% Anhänge
%% TODO: Inhalte festlegen, z. B. Screenshots, längere Code-Ausschnitte, Link zum Repository

\appsection{Quellcode}{anh:quellcode}
% Link zum Repository und kurze Beschreibung der Ordnerstruktur
TODO

% Beispiel für einen Screenshot im Anhang:
% \appsection{Screenshots der Benutzeroberfläche}{anh:screenshots}
% \begin{figure}[H]
% 	\centering
% 	\includegraphics[width=0.9\textwidth]{images/screenshot_dashboard.png}
% 	\caption{Dashboard zur Auswertung}
% 	\label{fig:dashboard}
% \end{figure}

	% \end{appendix}
    
	\newpage
\printbibliography[heading=bibintoc, title={\langLiteraturverzeichnis}]

    \ifcase\myKiDirectory
        \or
            \newpage
\section*{Verzeichnis der verwendeten KI-Hilfsmittel}
\markboth{Verzeichnis der verwendeten KI-Hilfsmittel}{Verzeichnis der verwendeten KI-Hilfsmittel}
\addcontentsline{toc}{section}{Verzeichnis der verwendeten KI-Hilfsmittel}

\begin{tabularx}{\textwidth}{|l|X|X|X|}
    \hline
    \textbf{KI-gestütztes Werkzeug} & \textbf{Prompts} & \textbf{Einsatzform} & \textbf{Seiten-/Bereichsangabe} \\
    \hline
    Claude (Anthropic) & Erstellung einer Gliederung für die Arbeit & Vorschlag für Gliederung & Gliederung \\
    \hline
\end{tabularx}

\newpage

        \else
    \fi


    \ifcase\myDecOfIndependence
        \or
            \thispagestyle{empty}

\begin{comment}
    \addcontentsline{toc}{section}{Selbständigkeitserklärung}
\end{comment}

\begin{center}
	\vspace*{2cm}
	\Huge\bf Selbständigkeitserklärung\\
	\vspace*{3cm}
	\large\rm\singlespacing 
	% Ich versichere hiermit, dass ich meine\\ 
    % \ifcase\myType Projektarbeit 
    % \or Bachelorarbeit 
    % \or \myTypeOfWork
    % \else\fi ~mit dem Thema\\
	% \vspace*{2cm}
	% \Large\bf\myTopic\\
	% \Large\rm\mySubTopic\\
	% \vspace*{2cm}
	% \large\rm
	% \singlespacing 
	% selbständig verfasst und keine anderen als die angegebenen\\Quellen und Hilfsmittel benutzt habe. 
	Ich versichere hiermit, dass ich die vorliegende Arbeit selbstständig verfasst und keine anderen als die angegebenen Quellen und Hilfsmittel verwendet habe und diese Arbeit bei keiner anderen Prüfung mit gleichem oder vergleichbarem Inhalt vorgelegt habe und diese bislang nicht veröffentlicht wurde. Des Weiteren versichere ich, dass die eingereichte elektronische Fassung mit der gedruckten Ausfertigung übereinstimmt.
	\vfill
	\begin{tabularx}{\textwidth}{l@{\extracolsep\fill}r}
            \myPlace, \myDate & \\
            \rule{7cm}{0.3mm} & \rule{7.55cm}{0.3mm}
    \end{tabularx}
    \begin{tabularx}{\textwidth}{*{2}{>{\arraybackslash}X}}
        Ort, Datum & Unterschrift 
    \end{tabularx}
\end{center}

            \thispagestyle{empty}
        \else
    \fi
 
\end{document}